\documentclass[11pt]{amsart}

\usepackage[utf8]{inputenc}
\usepackage[T1]{fontenc}
\usepackage{lmodern}
\usepackage[protrusion=true,expansion=false]{microtype}
\usepackage{amsmath,amssymb,amsthm}
\usepackage{mathtools}
\usepackage{booktabs}
\usepackage{tabularray}
\usepackage{enumitem}
\usepackage{xcolor}
\usepackage[
  colorlinks=true,
  linkcolor=black,
  citecolor=black,
  urlcolor=black
]{hyperref}

\hfuzz=2pt
\hbadness=5000
\vbadness=5000
\emergencystretch=2em
\raggedbottom
\allowdisplaybreaks[2]

\newtheorem{theorem}{Theorem}[section]
\newtheorem{proposition}[theorem]{Proposition}
\newtheorem{lemma}[theorem]{Lemma}
\newtheorem{corollary}[theorem]{Corollary}
\theoremstyle{definition}
\newtheorem{definition}[theorem]{Definition}
\theoremstyle{remark}
\newtheorem{remark}[theorem]{Remark}

\DeclareMathOperator{\MF}{MF}
\DeclareMathOperator{\Sing}{Sing}
\DeclareMathOperator{\CritLoc}{CritLoc}
\DeclareMathOperator{\Tot}{Tot}
\DeclareMathOperator{\Kar}{Kar}
\newcommand{\Db}{\mathrm D^{b}}
\newcommand{\Dtwo}{\mathrm{D}^{\mathrm{coh}}_{\mathbb Z/2}}
\newcommand{\Cliff}{\operatorname{Cl}}
\newcommand{\MorseCat}[1]{\mathsf{Q}_{#1}}
\newcommand{\MBcat}[1]{\mathsf{MB}(#1)}
\newcommand{\Aper}[2]{\mathsf{A}^{\mathrm{per}}_{#2}(#1)}
\newcommand{\QAper}[2]{\mathsf{QA}^{\mathrm{per}}_{#2}(#1)}
\newcommand{\QAtwist}[3]{\mathsf{QA}^{\mathrm{per}}_{#2}(#1;[#3])}
\newcommand{\Dper}[3]{\mathsf{D}^{\mathrm{per}}_{#2}(#1;[#3])}
\newcommand{\FormalCat}[1]{\mathsf{F}(#1)}
\newcommand{\CcompB}{\mathcal C^{\mathrm{comp}}_{B_4,0}}
\newcommand{\CcompV}{\mathcal C^{\mathrm{comp}}_{V_{18},0}}
\newcommand{\CrawB}{\mathcal C^{\mathrm{raw}}_{B_4,0}}
\newcommand{\CrawV}{\mathcal C^{\mathrm{raw}}_{V_{18},0}}
\newcommand{\CcompBtwo}{\mathcal C^{\mathrm{comp}}_{B_2,0}}
\newcommand{\CcompVten}{\mathcal C^{\mathrm{comp}}_{V_{10},0}}
\newcommand{\CcompBthree}{\mathcal C^{\mathrm{comp}}_{B_3,0}}
\newcommand{\CcompVfourteen}{\mathcal C^{\mathrm{comp}}_{V_{14},0}}
\newcommand{\CcompBone}{\mathcal C^{\mathrm{comp}}_{B_1,0}}
\newcommand{\CcompVsix}{\mathcal C^{\mathrm{comp}}_{V_6,0}}

\begin{document}

\title[Compactified gluing for $B_d$ and $V_{4d+2}$]{Compactified Gluing for the $B_d$ and $V_{4d+2}$ Landau--Ginzburg Models}
\author{Victor Ibrahim Santos El Adji}
\date{}

\begin{abstract}
We compare projective compactifications of the standard Landau--Ginzburg
models attached to the five rank-one pairs $(B_d,V_{4d+2})$, $1\leq d\leq5$.
For the genus-two pair $(B_4,V_{18})$, we carry out a complete chartwise
calculation: the four standard affine charts of $\mathbb P^3$, together with
the dense torus as a fifth computational chart, determine the central
singular supports, the noncentral ordinary double points, and the overlap
identifications.  The corresponding compactified central categories are
defined as explicit Čech totalizations of formal matrix-factorization
categories.  For all five pairs, we compare the isolated compactification
contributions and the numbers of irreducible components of the resolved
central fibers.  For the pair $(B_2,V_{10})$, we additionally compute the
source-torus critical strata, the raw central supports of convenient
projective pencils, and their two noncentral ordinary double points.  For
$(B_3,V_{14})$, we perform the analogous calculation and exhibit explicitly
the second $V_{14}$ node created by resolving a base curve.  On the
index-two side, we give compactification calculations for every $B_d$,
$1\leq d\leq5$: the quartic pencils, their isolated ordinary double points,
the base-locus singularities relevant to the resolution, and the central
component counts.  We also calculate the two remaining extreme cases on the
prime side directly from their compactified pencils.  The $V_6$ pencil has
one noncentral ordinary double point, while the $V_{22}$ pencil has three
ordinary double points away from the base locus and a fourth one created by
blowing up a boundary base line.  No additional source-torus calculation is
needed for these four extreme models.  The isolated and component-count
invariants agree for $d=2,3,4,5$ and differ for $d=1$.  We distinguish these numerical equalities
from an equivalence of the full compactified categories, which additionally
requires compatible global gluing data or compactified Homological Mirror
Symmetry.
\end{abstract}
\maketitle

\section{Scope, the five pairs, and the computational charts}

The rank-one families considered in this note form the five pairs
\[
(B_d,V_{4d+2}),
\qquad d=1,2,3,4,5.
\]
The uniform comparison records isolated compactification contributions and
central-component counts for every pair.  The complete chartwise resolution
is carried out for $(B_4,V_{18})$.  For $(B_2,V_{10})$ we give an explicit
critical-point and raw central-support calculation, while retaining the
standard resolution data for the component count.  For $(B_3,V_{14})$ we
also compute the source and compactifiable representatives, including the
base-curve blowup producing the second $V_{14}$ node.  For $B_1$ and $B_5$
we work directly with the compactified quartic pencils and do not repeat a
source-torus census.  We use the same compactified approach for $V_6$ and
$V_{22}$.  The $V_6$ pencil is naturally anticanonical in
$\mathbb P^1\times\mathbb P^2$; the other explicit projective pencils below
lie in $\mathbb P^3_{[X:Y:Z:T]}$.  For the latter, the projective cover has four
standard affine charts
\[
U_X,\qquad U_Y,\qquad U_Z,\qquad U_T.
\]
We add
\[
U_\circ=U_T\cap\{XYZ\neq0\}\simeq(\mathbb C^*)^3
\]
as a fifth, redundant, computational chart.  This chart records which
singularities occur before compactification.  Since $U_\circ\subset U_T$, a
point found in both charts is not a second singularity: the overlap maps must
identify the corresponding local categories.

The calculations in this note are reproduced by
\begin{center}
\texttt{python computations/rank\_one\_five\_chart\_gluing.py}.
\end{center}

There are three logically distinct objects:
\begin{enumerate}
\item the source Laurent polynomial used in the open-torus census;
\item a standard Givental-type Laurent representative that admits a convenient
quartic compactification;
\item the resolved Calabi--Yau Landau--Ginzburg model obtained after resolving
the base locus of the quartic pencil.
\end{enumerate}
Birational mutations can move critical strata between the torus and the
boundary.  We therefore do not identify the category of the source Laurent
polynomial with the compactified central category without an additional
argument.  Likewise, agreement of component or node counts for a general
pair does not determine its full exceptional-chart gluing.

\section{Reference survey and logical dependency map}
\label{sec:reference_survey}

This section records the role of each group of references.  The distinction is
important because the literature supplies the Laurent models, compactification
framework, and categorical tools, but it does not contain the five-chart
calculation or the claimed Čech presentations in exactly the form used below.

\subsection{Laurent models, periods, and compactifications}

The period tables and the classification of three-dimensional Minkowski
polynomials are developed in \cite{CCGK,Minkowski}.  For Picard-rank-one Fano
threefolds, explicit Laurent-polynomial weak Landau--Ginzburg models and open
Calabi--Yau compactifications are given case by case in
\cite{PrzyjalkowskiWeak}.  The formulas used below for $B_4$ and $V_{18}$ are
the standard toric/Givental representatives appearing in this circle of
constructions; the survey \cite{PrzyjalkowskiToric} provides a convenient
table of the rank-one models.  The compactifiable $B_2$ and $V_{10}$
representatives and their noncentral ordinary double points are also recorded
in \cite{Przyjalkowski2026}; the same source records the compactifiable
$B_3$ and $V_{14}$ pencils and the two isolated contributions of each.
For $V_6$ we use the Givental complete-intersection compactification, while
for $V_{22}$ we homogenize the Laurent polynomial of
\cite{PrzyjalkowskiWeak} and compare its base-locus resolution with
\cite{Przyjalkowski2026}.

The existence and geometry of Calabi--Yau compactifications, including the
behavior of the fiber over infinity, are treated systematically in
\cite{PrzyjalkowskiCY}.  The component counts and the explicit incidence
descriptions of the resolved central fibers of $B_4$ and $V_{18}$ are given in
\cite{PrzyjalkowskiHodge}.  The recent calculation
\cite{Przyjalkowski2026} concerns fibers with isolated singularities and is a
useful comparison for the two noncentral nodes, but it is not the primary
source for the central-fiber gluing carried out here.  In particular, the
homogeneous quartics, their affine chart equations, their Jacobian ideals, and
their overlap identifications below are direct calculations in this note.
For $d=2,3$, the cited sources provide the standard resolution and node
census, while the torus equations, raw central supports, Hessians, and the
displayed local blowup calculations are carried out directly in
Sections~\ref{sec:B2_V10_calculation} and
\ref{sec:B3_V14_calculation}.

\subsection{Mutation and the choice of a Laurent representative}

Mutation preserves the classical period and relates many Laurent mirrors of
the same Fano family \cite{Minkowski}.  For Givental/Hori--Vafa models coming
from different nef partitions, mutation equivalence is established in
\cite{PrinceMutation}.  These results justify comparing different Laurent
representatives at the level of periods and birational Landau--Ginzburg data.
They do not, without further hypotheses, identify their critical loci on a
fixed torus chart or their matrix-factorization categories after a chosen
compactification.  This is why the source Laurent polynomial, the convenient
quartic representative, and the resolved pencil are kept separate throughout
the paper.

\subsection{Matrix factorizations and formal localization}

Matrix factorizations originate in Eisenbud's description of maximal
Cohen--Macaulay modules over hypersurfaces \cite{Eisenbud}.  Orlov identifies
Landau--Ginzburg B-branes with categories of singularities in the affine case
\cite{OrlovDbranes} and develops the nonaffine theory in
\cite{OrlovNonaffine}.  Compact generation and dg models for local
matrix-factorization categories are studied in \cite{Dyckerhoff}.

The passage to formal neighborhoods and idempotent completion is supported by
Orlov's formal-completion theorem \cite{OrlovFormal}: after Karoubi completion,
the singularity category is determined by the formal neighborhood of the
singular locus.  The dg and support-sensitive framework needed to compare
overlaps and totalizations is developed further in \cite{Preygel}.  In this
paper, however, equations \eqref{eq:B4_cech_totalization} and
\eqref{eq:V18_cech_totalization} are definitions of the two compactified
central categories.  A complete descent theorem identifying those
totalizations with a globally constructed resolved matrix-factorization
category must be applied in an enhanced dg or stable-infinity-categorical
setting.

The local elimination of a split quadratic pair uses Kn\"orrer periodicity
\cite{Knorrer}; useful global and derived forms appear in
\cite{OrlovEquivalences,Hirano}.  Blowups are especially relevant because the
Calabi--Yau models are obtained by resolving a base locus.  The
matrix-factorization blowup formula and its dg foundations are given by
Lunts--Schn\"urer \cite{LuntsSchnurer}.  Their theorem explains what must be
tracked on every exceptional chart; it does not by itself show that all
exceptional contributions glue to a genus-two curve in the two pencils at
hand.

\subsection{The two genus-two categories}

For $B_4$, the intersection of two quadrics and its pencil of quadrics are
governed by the even Clifford algebra.  The resulting semiorthogonal
decomposition and its genus-two double cover are established in
\cite{KuzQuadric}; see also \cite{AuelBernardaraBolognesi} for a broader
treatment of quadric fibrations, Clifford algebras, and rationality questions.

For $V_{18}$, the homological projective dual of the $G_2$ Grassmannian and the
double cover of the orthogonal line branched at six points are constructed in
\cite{KuzHyperplane}.  This gives the genus-two curve and the
semiorthogonal decomposition used in \eqref{eq:V18_fano_sod}.  Kuznetsov's
survey \cite{KuzFano} formulates the relation between the nontrivial components
of index-two and index-one Fano threefolds.  The later work
\cite{KuzShinder} proves a deformation statement for the corresponding
components and corrects the scope of the original conjectural correspondence.
It should not be read as an automatic equivalence between an arbitrary
$B_4$ and an arbitrary $V_{18}$.

Finally, Bondal--Orlov reconstruction \cite{BondalOrlov} applies to the
$\mathbb Z$-graded derived categories of the smooth genus-two curves.  Passing
from that statement to the two-periodic categories in this paper requires a
compatible graded lift.  Passing from the Fano-side decompositions to
$\CcompB$ and $\CcompV$ additionally requires compactified Homological Mirror
Symmetry or a complete global formal-Clifford calculation.

\subsection{Summary of evidentiary status}

The reference chain may be summarized as follows:
\begin{center}
\small
\begin{tblr}{
  width=0.97\textwidth,
  colspec={Q[l,wd=0.21\textwidth]Q[l,wd=0.28\textwidth]X[l]},
  row{1}={font=\bfseries},
  hline{1,2,Z}={0.4pt},
  rowsep=2pt,
}
Input & Principal references & What they support here \\
Laurent mirrors & \cite{CCGK,PrzyjalkowskiWeak,PrzyjalkowskiToric} &
Choice of the Laurent representatives and their period interpretation. \\
Compactification & \cite{PrzyjalkowskiHodge,PrzyjalkowskiCY,Przyjalkowski2026} &
Existence and geometry of compactified pencils, central-component incidence, and comparison with isolated singular fibers. \\
Mutation & \cite{Minkowski,PrinceMutation} &
Period-preserving birational comparison, but not categorical invariance of the chosen compactification. \\
Local categories & \cite{Eisenbud,OrlovDbranes,OrlovNonaffine,Dyckerhoff} &
Matrix factorizations and singularity categories on affine and nonaffine charts. \\
Formal gluing & \cite{OrlovFormal,Preygel,LuntsSchnurer} &
Formal-local invariance, enhanced categorical operations, and behavior under blowups. \\
Quadratic reduction & \cite{Knorrer,OrlovEquivalences,Hirano} &
Removal of nondegenerate quadratic directions and the local curve-type model. \\
Fano-side curves & \cite{KuzQuadric,KuzHyperplane,KuzFano,KuzShinder} &
The two genus-two categories and the precise limits of their correspondence. \\
\end{tblr}
\end{center}

Thus the literature supports the inputs and the local categorical mechanisms.
The five-chart Jacobian calculations, the overlap audit, and the assertion that
no chartwise copy should be counted twice are the computational contribution of
this note.  The equivalence $\CcompB\simeq\CcompV$ remains conditional for the
reasons stated in Section~\ref{sec:reference_survey} and in the final section.

\section{The \texorpdfstring{$B_4$}{B4} pencil}

\subsection{The dense torus and projective equation}

For the chart audit we use
\begin{equation}
\label{eq:B4_torus}
f_{B_4}(x,y,z)=\frac{(x+1)^2(y+1)^2}{xyz}+z.
\end{equation}
This is a standard Givental-type representative of the $B_4$ mirror
\cite{PrzyjalkowskiWeak,PrzyjalkowskiToric}.
The fiber $f_{B_4}=\lambda$ has projective closure
\begin{equation}
\label{eq:B4_pencil}
F_{B_4,\lambda}
=(X+T)^2(Y+T)^2+XYZ^2-\lambda XYZT=0.
\end{equation}
The logarithmic critical equations of \eqref{eq:B4_torus} have the two
solutions
\[
(x,y,z)=(1,1,4),\qquad (1,1,-4),
\]
with critical values $8$ and $-8$.  Their logarithmic Hessian determinants
are $32$ and $-32$, respectively.

The corresponding projective points are
\[
p_+=[1:1:4:1],\qquad p_-=[1:1:-4:1].
\]
Both points have all homogeneous coordinates nonzero, so each appears in all
five computations.

\subsection{The four projective charts}

The point coordinates and affine Hessian determinants are
\[
\begin{array}{c|c|c|c}
\text{chart}&p_+&p_-&\det(\operatorname{Hess})\\ \hline
U_\circ,\ U_T&(1,1,4)&(1,1,-4)&128\\
U_X&(1,4,1)&(1,-4,1)&128\\
U_Y&(1,4,1)&(1,-4,1)&128\\
U_Z&(1/4,1/4,1/4)&(-1/4,-1/4,-1/4)&1/2.
\end{array}
\]
Thus the Čech overlap maps identify five local copies of $p_+$ with one
global point, and likewise for $p_-$.  The noncentral part contributes exactly
\[
\MorseCat{3}\oplus\MorseCat{3}.
\]
Here and below, ``noncentral part'' means the singularities away from the
unresolved base locus.

At $\lambda=0$, the singular support of the raw quartic is
\begin{equation}
\label{eq:B4_raw_support}
\begin{gathered}
L_X\cup L_Y\cup\{p_\infty\},\\
L_X=V(X+T,Z),\qquad L_Y=V(Y+T,Z),\qquad
p_\infty=[0:0:1:0].
\end{gathered}
\end{equation}
Write $G_i=F_{B_4,0}|_{U_i}$ and let
$J(G_i)=(G_i,\partial G_i)$ be its hypersurface Jacobian ideal.  With the
affine coordinates indicated by the formulas, the five equations are
\begin{align*}
G_\circ&=G_T|_{xyz\neq0},\\
G_T(x,y,z)&=(x+1)^2(y+1)^2+xyz^2,\\
G_X(y,z,t)&=(1+t)^2(y+t)^2+yz^2,\\
G_Y(x,z,t)&=(x+t)^2(1+t)^2+xz^2,\\
G_Z(x,y,t)&=(x+t)^2(y+t)^2+xy.
\end{align*}
Direct differentiation gives the following reduced Jacobian ideals:
\begin{align}
\sqrt{J(G_\circ)}&=(1),
\label{eq:B4_jacobian_torus}\\
\sqrt{J(G_T)}
&=(x+1,z)\cap(y+1,z),
\label{eq:B4_jacobian_T}\\
\sqrt{J(G_X)}
&=(t+1,z)\cap(y+t,z),
\label{eq:B4_jacobian_X}\\
\sqrt{J(G_Y)}
&=(x+t,z)\cap(t+1,z),
\label{eq:B4_jacobian_Y}\\
\sqrt{J(G_Z)}&=(x,y,t).
\label{eq:B4_jacobian_Z}
\end{align}
Thus the support is distributed among the charts as follows.
\[
\begin{array}{c|c}
\text{chart}&\text{central singular support}\\ \hline
U_\circ&\varnothing\\
U_T&\{x=-1,z=0\}\cup\{y=-1,z=0\}\\
U_X&\{t=-1,z=0\}\cup\{y+t=0,z=0\}\\
U_Y&\{x+t=0,z=0\}\cup\{t=-1,z=0\}\\
U_Z&\{(x,y,t)=(0,0,0)\}.
\end{array}
\]
The lines meet at
\[
q=[-1:-1:0:1].
\]
In $U_T$, put $X=-1+u$, $Y=-1+v$, and $Z=w$.  The central equation becomes
exactly
\begin{equation}
\label{eq:B4_crossing}
u^2v^2+(-1+u)(-1+v)w^2.
\end{equation}
Since the coefficient of $w^2$ is a unit near $q$, the formal splitting lemma
reduces \eqref{eq:B4_crossing} to
\[
w^2+(uv)^2.
\]
This is not a $D_4$ normal form.  Indeed,
\[
\CritLoc\bigl(w^2+u^2v^2\bigr)
=V(w,u)\cup V(w,v)
\]
is the union of two formal curves, whereas the article's relative $D_4$
model $p(p^2-\delta v^2)$ has isolated transverse critical locus
$V(p,v)$.  Therefore the $D_4$ notation $\Dper{Z}{4}{\delta}$ does not apply
at $q$.  More precisely, this is a \emph{nonisolated tangential
double-crossing singularity}:
\[
w^2+u^2v^2=(w+\mathrm iuv)(w-\mathrm iuv).
\]
The two factors define smooth formal surfaces.  At a generic point of either
critical curve they meet normally and the transverse singularity is $A_1$;
at $q$ their tangent planes coincide.  Its Jacobian algebra
\[
\mathbb C[[u,v,w]]/(w,uv^2,u^2v)
\]
is infinite-dimensional, so the germ is not an isolated ADE singularity.
The two line strata therefore form one connected gluing contribution; they
cannot be written as a direct sum of two Morse--Bott categories.

There are four further special points on the two lines:
\[
\begin{aligned}
a_X&=[-1:0:0:1],& b_X&=[0:1:0:0] &&\text{on }L_X,\\
a_Y&=[0:-1:0:1],& b_Y&=[1:0:0:0] &&\text{on }L_Y.
\end{aligned}
\]
They are distinct from $q$.  At $a_X$, use the $U_T$ coordinates
\[
X=-1+u,\qquad Y=s,\qquad Z=w.
\]
Then
\[
F_{B_4,0}=u^2(1+s)^2+(-1+u)s w^2
\overset{c}{\sim}U^2-sW^2.
\]
This is the nonisolated pinch-point, or $D_\infty$, normal form.  At $b_X$,
use $U_Y$ and set $T=s$, $X=u-s$, and $Z=w$; the equation is
\[
u^2(1+s)^2+(u-s)w^2
\overset{c}{\sim}U^2-sW^2.
\]
The points $a_Y$ and $b_Y$ have the same formal type by symmetry.  Thus the
transverse $A_1$ form along either line degenerates at two pinch points, in
addition to the tangential double crossing at $q$.

In the $U_Z$ chart around $p_\infty$, the central equation is
\[
xy+(x+t)^2(y+t)^2.
\]
Its quadratic part is $xy$.  Eliminating the two nondegenerate variables
leaves a residual series beginning with $t^4$, so the raw quartic has an
$A_3$ singularity at $p_\infty$; equivalently, its local Jacobian ideal is
$(x,y,t^3)$.  This point is singular in the base-locus
model and is modified by the simultaneous resolution.  Consequently its raw
$A_3$ category must not be added as an independent summand of the resolved
Landau--Ginzburg model.

\subsection{Decomposition on each of the five raw charts}

For a chart $U_i$, let $\mathcal R^{B_4}_{i,0}$ denote the
idempotent-completed formal matrix-factorization category of the raw central
equation in that chart.  For $i\in\{T,X,Y\}$, the notation
$\FormalCat{(L_X\cup L_Y)\cap U_i}$ denotes the central formal category
supported on
$(L_X\cup L_Y)\cap U_i$.  Away from
\[
\Sigma_{B_4}=\{q,a_X,b_X,a_Y,b_Y\}
\]
the transverse Hessian is nondegenerate, hence locally the contribution is
Morse--Bott.  At the four points $a_X,b_X,a_Y,b_Y$ the formal type is
$D_\infty$, while at $q$ it is $w^2+(uv)^2$.  These special germs remain
joined to the open Morse--Bott strata inside
$\FormalCat{(L_X\cup L_Y)\cap U_i}$.

The raw central categories are
\begin{align}
\mathcal R^{B_4}_{\circ,0}
&\simeq 0,
\label{eq:B4_chart_decomposition_torus}\\
\mathcal R^{B_4}_{T,0}
&\simeq \FormalCat{(L_X\cup L_Y)\cap U_T},\\
\mathcal R^{B_4}_{X,0}
&\simeq \FormalCat{(L_X\cup L_Y)\cap U_X},\\
\mathcal R^{B_4}_{Y,0}
&\simeq \FormalCat{(L_X\cup L_Y)\cap U_Y},\\
\mathcal R^{B_4}_{Z,0}
&\simeq \Aper{\mathrm{pt}}{3}.
\label{eq:B4_chart_decomposition_Z}
\end{align}
More explicitly, equations \eqref{eq:B4_jacobian_torus}--
\eqref{eq:B4_jacobian_Z} give the dictionary
\begin{center}
\small
\begin{tblr}{
  width=\textwidth,
  colspec={Q[c,wd=0.08\textwidth]X[l]Q[l,wd=0.31\textwidth]},
  row{1}={font=\bfseries},
  hline{1,2,Z}={0.4pt},
  rowsep=2pt,
}
chart & formal type on the support & category \\
$U_\circ$ & empty & $0$ \\
$U_T$ & $w^2+(uv)^2$ at $q$; $D_\infty$ at $a_X,a_Y$; $A_1$ elsewhere
& $\FormalCat{(L_X\cup L_Y)\cap U_T}$ \\
$U_X$ & $w^2+(uv)^2$ at $q$; $D_\infty$ at $a_X,b_Y$; $A_1$ elsewhere
& $\FormalCat{(L_X\cup L_Y)\cap U_X}$ \\
$U_Y$ & $w^2+(uv)^2$ at $q$; $D_\infty$ at $b_X,a_Y$; $A_1$ elsewhere
& $\FormalCat{(L_X\cup L_Y)\cap U_Y}$ \\
$U_Z$ & $xy+t^4$ & $\Aper{\mathrm{pt}}{3}$ \\
\end{tblr}
\end{center}
Here $xy+t^4$ denotes the formal contact type after eliminating higher terms;
it is precisely the convention $\Aper{\mathrm{pt}}{3}$ used in the main
article.  The three occurrences of the line-crossing category are restrictions
of one global connected category, not three direct summands.
In addition, the part supported away from the base locus is
$\MorseCat{3}^{\oplus2}$ in every chart, because both noncentral nodes have all
homogeneous coordinates nonzero.  These are five restrictions of two global
points, not ten global summands.  Likewise, on pairwise and higher overlaps the restrictions of
the three displayed $\FormalCat{(L_X\cup L_Y)\cap U_i}$ categories are
identified.  The $A_3$ point occurs only in $U_Z$
and is disjoint from the line crossing in the raw central fiber.

Because the raw pencil still has an unresolved base locus, the expressions
above are not asserted to be the full matrix-factorization categories of all
raw fibers.  They record exactly the central fiber and the two verified
singularities away from the base locus.  The resolved Landau--Ginzburg map is
needed before taking a genuine direct sum over all critical values.

\subsection{Raw and resolved central categories}

Let $\FormalCat{L_X\cup L_Y}$ denote the formal matrix-factorization category
supported on the connected crossing $L_X\cup L_Y$.  Before resolving the
base locus, disjointness of $p_\infty$ gives
\begin{equation}
\label{eq:B4_raw_category}
\CrawB
\simeq
\FormalCat{L_X\cup L_Y}
\oplus \Aper{\mathrm{pt}}{3}.
\end{equation}
Equation \eqref{eq:B4_raw_category} is a statement about the raw quartic, not
the resolved compactification.

This formula can be refined by supported localization.  Put
\[
L_X^\circ=L_X\setminus\Sigma_{B_4},
\qquad
L_Y^\circ=L_Y\setminus\Sigma_{B_4},
\]
and define the category retaining all five special germs by
\[
\mathcal P^{\mathrm{raw}}_{B_4}
:=
\overline{\MF^{\Sigma_{B_4}}_0\!\left(
\widehat{\mathbb P^3}_{L_X\cup L_Y},\widehat F_{B_4,0}
\right)}.
\]
Its completed local germs consist of four $D_\infty$ pinch points and the
tangential double crossing $w^2+(uv)^2$ at $q$.  Removing
$\Sigma_{B_4}$ leaves the disjoint union
$L_X^\circ\amalg L_Y^\circ$, and supported localization gives
\begin{equation}
\label{eq:B4_raw_localization}
\mathcal P^{\mathrm{raw}}_{B_4}
\longrightarrow
\FormalCat{L_X\cup L_Y}
\longrightarrow
\MBcat{L_X^\circ}\oplus\MBcat{L_Y^\circ}.
\end{equation}
The notation $\MBcat{L_i^\circ}$ retains the actual transverse quadratic
bundle and its Clifford data.  Sequence \eqref{eq:B4_raw_localization} is not
asserted to split.  Since $p_\infty$ is disjoint from both lines, its
$\Aper{\mathrm{pt}}{3}$ contribution is the genuine direct summand appearing
in \eqref{eq:B4_raw_category}.

Let $\pi_{B_4}:\widetilde Y_{B_4}\to Y^{\mathrm{raw}}_{B_4}$ be the
standard minimal resolution.  Its central fiber has three surface components
\[
D_1\cup D_2\cup D_3.
\]
Their pairwise intersections
\[
R_{12}=D_1\cap D_2,
\qquad
R_{13}=D_1\cap D_3,
\qquad
R_{23}=D_2\cap D_3
\]
are rational curves, and the three curves meet at two points
$e_+$ and $e_-$ \cite{PrzyjalkowskiHodge}.  Thus the resolved central
singular support is the connected curve
\[
\Gamma_{B_4}=R_{12}\cup R_{13}\cup R_{23},
\qquad
R_{12}\cap R_{13}\cap R_{23}=\{e_+,e_-\}.
\]
Put
\[
R_{ij}^{\circ}=R_{ij}\setminus\{e_+,e_-\}.
\]
At a point of $R_{ij}^{\circ}$ the resolved map has transverse equation
$uv$, so that the article's notation for this open contribution is
$\MBcat{R_{ij}^{\circ}}$.  For $\epsilon\in\{+,-\}$, let
\[
\mathcal P_{B_4,\epsilon}
:=
\overline{\MF^{e_\epsilon}_0
\bigl(\widehat{\widetilde Y}_{B_4,e_\epsilon},
\widehat w_{B_4}\bigr)},
\qquad
\mathcal P_{B_4}:=
\mathcal P_{B_4,+}\oplus\mathcal P_{B_4,-}.
\]
This definition retains the complete formal equation at the two triple
points.  If a chosen final resolution chart is simple normal crossing, then
that equation is $uvw$ and
\[
\mathcal P_{B_4,\epsilon}
\simeq
\overline{\MF^t_0
\bigl(\operatorname{Spf}\mathbb C[[u,v,w]],uvw\bigr)}.
\]
The component-count statement alone does not replace this last local
normal-form check.
Supported localization gives the explicit nonsplit Verdier sequence
\begin{equation}
\label{eq:B4_resolved_localization}
\mathcal P_{B_4}
\longrightarrow
\CcompB
\longrightarrow
\MBcat{R_{12}^{\circ}}
\oplus\MBcat{R_{13}^{\circ}}
\oplus\MBcat{R_{23}^{\circ}}.
\end{equation}
In particular, the three rational-curve terms in
\eqref{eq:B4_resolved_localization} are not asserted to form a direct sum
decomposition of $\CcompB$; the two triple-point categories retain their
gluing.

Put $\widetilde U_i=\pi_{B_4}^{-1}(U_i)$.  For a nonempty
subset $I\subset\{\circ,X,Y,Z,T\}$, put
\[
\widetilde U_I=\bigcap_{i\in I}\widetilde U_i,
\qquad
\mathfrak M^B_I=
\MF^{\Gamma_{B_4}\cap\widetilde U_I}_{\mathrm{dg}}
\bigl(\widehat{\widetilde U_I}_{\Gamma_{B_4}\cap\widetilde U_I},
\widehat w_{B_4}\bigr).
\]
For every inclusion $I\subset J$, restriction along
$\widetilde U_J\hookrightarrow\widetilde U_I$ gives a dg functor
\begin{equation}
\label{eq:B4_cech_restriction}
\rho^B_{I,J}:\mathfrak M^B_I\longrightarrow\mathfrak M^B_J.
\end{equation}
Together with their canonical composition coherences, these functors form a
semicosimplicial dg category
\[
\check{\mathfrak M}^{B,p}
=
\prod_{i_0<\cdots<i_p}
\mathfrak M^B_{\{i_0,\ldots,i_p\}}.
\]
On the component indexed by $J=\{j_0<\cdots<j_{p+1}\}$, the $a$-th coface
is $\rho^B_{J\setminus\{j_a\},J}$.

\begin{definition}
\label{def:B4_comp_category}
The compactified central category is the idempotent-completed Čech
totalization of this enhanced diagram:
\begin{equation}
\label{eq:B4_cech_totalization}
\CcompB
:=
\Kar\!\left(
H^0\Tot\bigl(\check{\mathfrak M}^{B,\bullet}\bigr)
\right).
\end{equation}
\end{definition}

Here totalization means the homotopy limit in the dg, equivalently
stable-$\infty$-categorical, setting.  Equation
\eqref{eq:B4_cech_totalization} is a definition; it does not assert that the
outer terms of the Verdier sequence \eqref{eq:B4_resolved_localization}
determine its middle term.

The overlap identifications and the two noncentral Morse points give the
proved top-level decomposition
\begin{equation}
\label{eq:B4_final_decomposition}
\overline{\MF(Y_{B_4},w_{B_4})}
\simeq
\CcompB\oplus\MorseCat{3}\oplus\MorseCat{3}.
\end{equation}

\section{The \texorpdfstring{$V_{18}$}{V18} pencil}

\subsection{The torus model and its two ordinary double points}

Use the standard representative
\begin{equation}
\label{eq:V18_torus}
f_{V_{18}}(x,y,z)
=\frac{(x+y+z)(x+y+z+xy+xz+yz+xyz)}{xyz}.
\end{equation}
The corresponding compactified pencil belongs to the standard toric
compactification framework of
\cite{PrzyjalkowskiWeak,PrzyjalkowskiCY,PrzyjalkowskiToric}.
Its projective pencil is
\begin{equation}
\label{eq:V18_pencil}
\begin{split}
F_{V_{18},\lambda}
={}&(X+Y+Z)
\bigl((X+Y+Z)T^2+(XY+XZ+YZ)T+XYZ\bigr)\\
&-\lambda XYZT=0.
\end{split}
\end{equation}
Put $3a^2=1$.  The two singular points outside the base locus are
\begin{equation}
\label{eq:V18_nodes}
p_a=[1:1:1:a],
\qquad
\lambda=18a+9=9\pm6\sqrt3.
\end{equation}
Their coordinates are
\[
\begin{array}{c|c}
\text{chart}&p_a\\ \hline
U_\circ,\ U_T&(a^{-1},a^{-1},a^{-1})\\
U_X&(1,1,a)\\
U_Y&(1,1,a)\\
U_Z&(1,1,a).
\end{array}
\]
The Hessian determinant is invertible modulo $3a^2-1$ in every chart.  Hence
the five local descriptions glue to two global ordinary double points and
contribute $\MorseCat{3}\oplus\MorseCat{3}$.

\subsection{The connected central cubic}

At $\lambda=0$, set
\[
A=X+Y+Z,
\qquad
B=(X+Y+Z)T^2+(XY+XZ+YZ)T+XYZ.
\]
The central quartic is $AB=0$.  Its two components meet along
\begin{equation}
\label{eq:V18_cubic}
C_0=V(A,B),
\qquad
T(X^2+XY+Y^2)+XY(X+Y)=0
\quad\text{on }A=0.
\end{equation}
Since $A$ is linear and smooth, a point of $AB=0$ is singular either when
$A=B=0$ or when it is a singular point of the cubic surface $B=0$.  Direct
differentiation gives
\[
\Sing(V(B))=\{r_X,r_Y,r_Z\},
\]
so the complete singular support of the raw central quartic is
\[
\Sing(V(AB))=C_0\sqcup\{r_X,r_Y,r_Z\}.
\]
The chart equations are
\[
\begin{array}{c|c}
\text{chart}&C_0\\ \hline
U_\circ&x+y+z=0,\quad xy+xz+yz+xyz=0,\quad xyz\neq0\\
U_T&x+y+z=0,\quad xy+xz+yz+xyz=0\\
U_X&1+y+z=0,\quad t(y+z+yz)+yz=0\\
U_Y&x+1+z=0,\quad t(x+z+xz)+xz=0\\
U_Z&x+y+1=0,\quad t(x+y+xy)+xy=0.
\end{array}
\]
For complete explicitness, put
\begin{align*}
A_T&=x+y+z,
&B_T&=A_T+xy+xz+yz+xyz,\\
A_X&=1+y+z,
&B_X&=A_X\,t^2+(y+z+yz)t+yz,\\
A_Y&=x+1+z,
&B_Y&=A_Y\,t^2+(x+z+xz)t+xz,\\
A_Z&=x+y+1,
&B_Z&=A_Z\,t^2+(x+y+xy)t+xy.
\end{align*}
Then $G_i=A_iB_i$ on each projective chart, while
$G_\circ=G_T|_{xyz\neq0}$.  If $J(G_i)=(G_i,\partial G_i)$, direct
differentiation gives
\begin{align}
\sqrt{J(G_\circ)}&=(A_T,B_T)\mathcal O(U_\circ),
\label{eq:V18_jacobian_torus}\\
\sqrt{J(G_T)}&=(A_T,B_T),
\label{eq:V18_jacobian_T}\\
\sqrt{J(G_X)}&=(A_X,B_X)\cap(y,z,t),
\label{eq:V18_jacobian_X}\\
\sqrt{J(G_Y)}&=(A_Y,B_Y)\cap(x,z,t),
\label{eq:V18_jacobian_Y}\\
\sqrt{J(G_Z)}&=(A_Z,B_Z)\cap(x,y,t).
\label{eq:V18_jacobian_Z}
\end{align}
The first factor in each intersection is the corresponding piece of $C_0$;
the maximal ideals in the last three lines are respectively the points
$r_X,r_Y,r_Z$.
The plane cubic \eqref{eq:V18_cubic} is irreducible and has one ordinary node
at
\[
p_0=[0:0:0:1].
\]
In $U_T$ the quadratic term at $p_0$ is $x^2+xy+y^2$, whose Hessian is
nondegenerate.  Away from $p_0$, the gradients of $A$ and $B$ are independent,
so the formal local equation is $uv$.  At $p_0$ the two equations are tangent,
and the node cannot be detached as an independent direct summand.

The seven points
\[
\begin{gathered}
[1:0:0:0],\quad[0:1:0:0],\quad[0:0:1:0],\\
[1:-1:0:0],\quad[1:0:-1:0],\quad[0:1:-1:0],\\
[0:0:0:1]
\end{gathered}
\]
are special points of the singular base-locus model.  Their roles are not the
same.  The three coordinate points
\[
r_X=[1:0:0:0],\qquad
r_Y=[0:1:0:0],\qquad
r_Z=[0:0:1:0]
\]
are disjoint from $C_0$ and are isolated singularities of the raw central
quartic.  For example, in $U_X$ the local central equation at $r_X$ is
\[
(1+y+z)
\bigl(t^2(1+y+z)+t(y+z+yz)+yz\bigr).
\]
Its Hessian has rank two and kernel generated by
$(y,z,t)=(-s,-s,s)$.  Restriction to this kernel is
\[
s^3(2s-1),
\]
so the residual cubic is nonzero and $r_X$ has type $A_2$.  Cyclic symmetry
gives the same type at $r_Y$ and $r_Z$.  This is the central-fiber
calculation at $\lambda=0$.  For the general pencil member, the cubic
coefficient along the Hessian kernel is $-(1+\lambda)s^3$; at $\lambda=-1$
it vanishes and each of the three points has type $A_3$.  This degeneration
lies on the fixed base locus and is resolved in the simultaneous resolution;
it does not create an additional ordinary double point, consistently with
$\varepsilon(V_{18})=2$.

The three points
\[
[1:-1:0:0],\qquad[1:0:-1:0],\qquad[0:1:-1:0]
\]
are smooth points of $C_0$.  The remaining point $p_0=[0:0:0:1]$ is its node.
All seven points are affected by the resolution of the base locus, but only
$r_X,r_Y,r_Z$ give additional isolated summands for the raw quartic.  None of
these raw summands may be inserted unchanged into the resolved category
\cite{PrzyjalkowskiCY,PrzyjalkowskiToric}.

The incidence relations are therefore
\[
V(A)\cap V(B)=C_0,
\qquad
C_0\cap\{r_X,r_Y,r_Z\}=\varnothing,
\qquad
r_i\neq r_j\quad(i\neq j).
\]
Put $C_0^\circ=C_0\setminus\{p_0\}$ and define the category supported at the
node by
\[
\mathcal P^{\mathrm{raw}}_{V_{18},p_0}
:=
\overline{\MF^{p_0}_0\!\left(
\widehat{\mathbb P^3}_{C_0},\widehat F_{V_{18},0}
\right)}.
\]
Its local formal equation is $w^2+(uv)^2$.  Along $C_0^\circ$, the functions
$(A,B)$ are independent formal normal coordinates and the equation is exactly
$AB=uv$.  Supported localization consequently gives the nonsplit sequence
\begin{equation}
\label{eq:V18_raw_localization}
\mathcal P^{\mathrm{raw}}_{V_{18},p_0}
\longrightarrow
\FormalCat{C_0}
\longrightarrow
\MBcat{C_0^\circ}.
\end{equation}
The isolated points are disjoint from this connected support, and their local
matrix-factorization categories are, for $i\in\{X,Y,Z\}$,
\[
\mathcal A_i
:=
\overline{\MF^{r_i}_0(\widehat{U_i}_{r_i},\widehat F_{V_{18},0})}
\simeq\Aper{\mathrm{pt}}{2}.
\]
Consequently
$\mathcal A_X\oplus\mathcal A_Y\oplus\mathcal A_Z
\simeq\Aper{\mathrm{pt}}{2}^{\oplus3}$.

\subsection{Decomposition on each of the five raw charts}

Put $C_{0,i}=C_0\cap U_i$.  The opens
$C_{0,\circ},C_{0,X},C_{0,Y},C_{0,Z}$ are smooth, and the exact product
$AB$ has split normal form $uv$ along them.  Kn\"orrer periodicity therefore
identifies their central chart contributions with
$\MBcat{C_{0,i}}$.  The chart $U_T$ contains the node $p_0$, and
$\FormalCat{C_{0,T}}$ denotes the single formal category supported on all of
$C_{0,T}$.

To see its special local type, use $r=A=x+y+z$ as one coordinate and put
$z=r-x-y$.  Writing $B=r+h$ gives the exact square completion
\[
AB=r(r+h)=\left(r+\frac h2\right)^2-\left(\frac h2\right)^2.
\]
The restriction $h|_{r=0}$ has nondegenerate quadratic part
$-(x^2+xy+y^2)$.  The formal splitting and Morse lemmas therefore give the
same crossing type $w^2+(uv)^2$ at $p_0$.  Thus the two formal branches of the
node belong to one category $\FormalCat{C_{0,T}}$ rather than two direct
summands.  As in the $B_4$ calculation, this is a nonisolated tangential
double crossing: it is generically transverse $A_1$ along each branch, but at
$p_0$ the two smooth surface branches are tangent.  It is not the isolated
transverse model $p(p^2-\delta v^2)$ defining $\Dper{Z}{4}{\delta}$.

Let $\mathcal R^{V_{18}}_{i,0}$ denote the raw central category in $U_i$.
The chartwise singular supports and the preceding local normal forms give
\begin{align}
\mathcal R^{V_{18}}_{\circ,0}
&\simeq \MBcat{C_{0,\circ}},
\label{eq:V18_chart_decomposition_torus}\\
\mathcal R^{V_{18}}_{T,0}
&\simeq \FormalCat{C_{0,T}},\\
\mathcal R^{V_{18}}_{X,0}
&\simeq
\MBcat{C_{0,X}}
\oplus\Aper{\mathrm{pt}}{2},\\
\mathcal R^{V_{18}}_{Y,0}
&\simeq
\MBcat{C_{0,Y}}
\oplus\Aper{\mathrm{pt}}{2},\\
\mathcal R^{V_{18}}_{Z,0}
&\simeq
\MBcat{C_{0,Z}}
\oplus\Aper{\mathrm{pt}}{2}.
\label{eq:V18_chart_decomposition_Z}
\end{align}
Equivalently, the calculation may be read directly from the Jacobian ideals:
\[
\begin{array}{c|c|c}
\text{chart}&\text{formal type on the support}&\text{category}\\ \hline
U_\circ&uv\text{ along }C_{0,\circ}
&\MBcat{C_{0,\circ}}\\
U_T&uv\text{ off }p_0,\quad w^2+(uv)^2\text{ at }p_0
&\FormalCat{C_{0,T}}\\
U_X&uv\text{ along }C_{0,X},\quad uv+s^3\text{ at }r_X
&\MBcat{C_{0,X}}\oplus\Aper{\mathrm{pt}}{2}\\
U_Y&uv\text{ along }C_{0,Y},\quad uv+s^3\text{ at }r_Y
&\MBcat{C_{0,Y}}\oplus\Aper{\mathrm{pt}}{2}\\
U_Z&uv\text{ along }C_{0,Z},\quad uv+s^3\text{ at }r_Z
&\MBcat{C_{0,Z}}\oplus\Aper{\mathrm{pt}}{2}.
\end{array}
\]
The notation $uv+s^3$ records the three-variable suspension of the $A_2$
residual singularity, hence exactly the article's term
$\Aper{\mathrm{pt}}{2}$.
Here the three $A_2$ terms are genuinely different raw points: $r_X$ occurs
only in $U_X$, $r_Y$ only in $U_Y$, and $r_Z$ only in $U_Z$.  By contrast,
the curve pieces are identified by the overlap functors.  Separately, the part
supported away from the base locus is $\MorseCat{3}^{\oplus2}$ in every chart;
its five chartwise copies glue to the same two global nodes.

As for $B_4$, these formulas do not define the full categories of every fiber
of the unresolved rational pencil: base-locus singularities can persist for
other values of $\lambda$.  They give the complete decomposition of the raw
central fiber together with the separately verified non-base-locus nodes.

The notation $\FormalCat{C_0}$ denotes the connected global formal category
along the nodal cubic appearing in
\eqref{eq:V18_raw_localization}.  Before resolving the base locus, the central
category is consequently
\begin{equation}
\label{eq:V18_raw_category}
\CrawV
\simeq
\FormalCat{C_0}
\oplus
\Aper{\mathrm{pt}}{2}^{\oplus3}.
\end{equation}
This is a decomposition of the raw central quartic only.

Let $\pi_{V_{18}}:\widetilde Y_{V_{18}}\to
Y^{\mathrm{raw}}_{V_{18}}$ be the standard minimal resolution.  Its central
fiber has three surface components
\[
E_1\cup E_2\cup E_3.
\]
Each pair meets in a rational curve
\[
S_{12}=E_1\cap E_2,
\qquad
S_{13}=E_1\cap E_3,
\qquad
S_{23}=E_2\cap E_3,
\]
and all three curves meet at two points $d_+$ and $d_-$
\cite{PrzyjalkowskiHodge}.  Therefore
\[
\Gamma_{V_{18}}=S_{12}\cup S_{13}\cup S_{23}
\]
is the complete connected singular support of the resolved central fiber.
These two points lie over the node $p_0=[0:0:0:1]$ and record its two tangent
directions.  Indeed, in $U_T$ and on the plane $A=0$, the nodal cubic has
equation
\[
g(x,y)=x^2+xy+y^2+xy(x+y)=0.
\]
Blow up $(x,y)=(0,0)$ and use the chart
\[
x=\tau s,\qquad y=s.
\]
After removing the exceptional factor $s^2$, the strict transform is
\[
\tau^2+\tau+1+\tau(\tau+1)s=0.
\]
Hence its intersection with the exceptional divisor $s=0$ consists of
\[
d_+=\{s=0,\ \tau=\omega_+\},
\qquad
d_-=\{s=0,\ \tau=\omega_-\},
\]
where
\[
\omega_\pm=\frac{-1\pm\mathrm i\sqrt3}{2},
\qquad
\omega_\pm^2+\omega_\pm+1=0.
\]
Equivalently,
\[
\pi_{V_{18}}(d_+)=\pi_{V_{18}}(d_-)=p_0,
\qquad
[x:y]_{d_\pm}=[\omega_\pm:1].
\]
Put $S_{ij}^{\circ}=S_{ij}\setminus\{d_+,d_-\}$ and define
\[
\mathcal P_{V_{18},\epsilon}
:=
\overline{\MF^{d_\epsilon}_0
\bigl(\widehat{\widetilde Y}_{V_{18},d_\epsilon},
\widehat w_{V_{18}}\bigr)},
\qquad
\mathcal P_{V_{18}}
:=
\mathcal P_{V_{18},+}\oplus\mathcal P_{V_{18},-}.
\]
If the final resolution chart is simple normal crossing, then
$\mathcal P_{V_{18},\epsilon}$ is the local category of $uvw$; otherwise the
displayed definition retains the actual formal germ.  Supported localization
gives
\begin{equation}
\label{eq:V18_resolved_localization}
\mathcal P_{V_{18}}
\longrightarrow
\CcompV
\longrightarrow
\MBcat{S_{12}^{\circ}}
\oplus\MBcat{S_{13}^{\circ}}
\oplus\MBcat{S_{23}^{\circ}}.
\end{equation}
Again this is a gluing statement, not a direct-sum decomposition.

Put $\widetilde U_i=\pi_{V_{18}}^{-1}(U_i)$.
For each nonempty $I\subset\{\circ,X,Y,Z,T\}$, define
\[
\widetilde U_I=\bigcap_{i\in I}\widetilde U_i,
\qquad
\mathfrak M^V_I=
\MF^{\Gamma_{V_{18}}\cap\widetilde U_I}_{\mathrm{dg}}
\bigl(\widehat{\widetilde U_I}_{\Gamma_{V_{18}}\cap\widetilde U_I},
\widehat w_{V_{18}}\bigr).
\]
For $I\subset J$, restriction gives
\begin{equation}
\label{eq:V18_cech_restriction}
\rho^V_{I,J}:\mathfrak M^V_I\longrightarrow\mathfrak M^V_J.
\end{equation}
These functors and their composition coherences define
\[
\check{\mathfrak M}^{V,p}
=
\prod_{i_0<\cdots<i_p}
\mathfrak M^V_{\{i_0,\ldots,i_p\}},
\]
whose cofaces are $\rho^V_{J\setminus\{j_a\},J}$.

\begin{definition}
\label{def:V18_comp_category}
The compactified central category is
\begin{equation}
\label{eq:V18_cech_totalization}
\CcompV
:=
\Kar\!\left(
H^0\Tot\bigl(\check{\mathfrak M}^{V,\bullet}\bigr)
\right).
\end{equation}
\end{definition}

Again the totalization is taken before passing to the homotopy category and
idempotent completion.

Thus the proved decomposition is
\begin{equation}
\label{eq:V18_final_decomposition}
\overline{\MF(Y_{V_{18}},w_{V_{18}})}
\simeq
\CcompV\oplus\MorseCat{3}\oplus\MorseCat{3}.
\end{equation}

\section{The \texorpdfstring{$d=2$}{d=2} pair: \texorpdfstring{$B_2$ and $V_{10}$}{B2 and V10}}
\label{sec:B2_V10_calculation}

We now perform the corresponding critical-point calculation for the quartic
double solid $B_2$ and the Gushel--Mukai threefold $V_{10}$.  As in the
$d=4$ calculation, it is essential to distinguish the source Laurent
polynomials used in the torus census from the convenient Laurent
representatives used for projective compactification.

\subsection{The source Laurent polynomials on the torus}

For the source $B_2$ model, set
\begin{equation}
\label{eq:B2_source_factorization}
F_{B_2}=x(y+z+1)+\frac1{xyz},
\qquad
W_{B_2}=\frac{F_{B_2}^{,2}}{x}.
\end{equation}
The surface
\[
S_{B_2}=V(F_{B_2})\subset(\mathbb C^*)^3
\]
is smooth.  Along it the potential has a rank-one normal Hessian with
quadratic coefficient $1/x$, whose square class is $[x]$.  This class is
nontrivial: on the curve
\[
(x,y,z)=(t,-1,t^{-1})\subset S_{B_2}
\]
it restricts to $[t]$.  Saturation of the logarithmic critical ideal by
$F_{B_2}$ gives the two complementary points
\[
p_+=(1,1,1),
\qquad
p_-=(-1,1,1).
\]
Their critical values and logarithmic Hessian determinants are
\[
W_{B_2}(p_\pm)=\pm16,
\qquad
\det(\theta_i\theta_jW_{B_2})(p_\pm)=\pm512.
\]
Consequently,
\begin{equation}
\label{eq:B2_source_category}
\overline{\MF((\mathbb C^*)^3,W_{B_2})}
\simeq
\MBcat{S_{B_2};[x]}\oplus\MorseCat{3}^{\oplus2}.
\end{equation}

For the source $V_{10}$ polynomial of the torus census, put
\[
F_{10}=xyz^2+xyz+1,
\qquad
G_{10}=xyz^2+xyz+xz+yz+1.
\]
Its $21$ monomials group into the exact factorization
\begin{equation}
\label{eq:V10_source_factorization}
W_{V_{10}}+6
=
\frac{F_{10}^{,2}G_{10}}{x^2y^2z^3}.
\end{equation}
Thus the fiber at $-6$ contains the smooth critical surface
\[
S_{V_{10}}=V(F_{10}).
\]
The coefficient $G_{10}$ is a unit on
\[
S_{V_{10}}^\circ=S_{V_{10}}\setminus K,
\qquad
K=V(F_{10},G_{10}),
\]
and $K$ has the explicit presentation
\[
K=\{y=-x,\ x^2z(z+1)=1\}.
\]
Moreover,
\[
\det\!\left(
\frac{\partial(F_{10},G_{10})}{\partial(x,y)}
\right)\bigg|_K
=-2xz^2(z+1),
\]
which is a unit on $K$.  Hence $(F_{10},G_{10})$ are formal normal
coordinates along $K$, and the shifted potential is contact-equivalent there
to $u^2v$.  If $\mathcal P_{V_{10}}$ denotes the torsion formal category of
this crossing family and
\[
\mathcal C_{V_{10},-6}
:=
\overline{\MF^{S_{V_{10}}}
(\widehat{(\mathbb C^*)^3}_{S_{V_{10}}},\widehat W_{V_{10}})},
\]
supported localization gives
\begin{equation}
\label{eq:V10_source_localization}
\mathcal P_{V_{10}}
\longrightarrow
\mathcal C_{V_{10},-6}
\longrightarrow
\Aper{S_{V_{10}}^\circ}{1}.
\end{equation}
The remaining critical points are
\[
P_a=(4a-2,4a-2,a),
\qquad
4a^2-2a-1=0,
\]
and satisfy
\[
W_{V_{10}}(P_a)=40a+6,
\qquad
\det(\theta_i\theta_jW_{V_{10}})(P_a)=64(123a+38)\neq0.
\]
It follows that
\begin{equation}
\label{eq:V10_source_category}
\overline{\MF((\mathbb C^*)^3,W_{V_{10}})}
\simeq
\mathcal C_{V_{10},-6}\oplus\MorseCat{3}^{\oplus2}.
\end{equation}
Equations \eqref{eq:B2_source_category} and
\eqref{eq:V10_source_category} already show that equality of the two source
open-torus categories does not follow from their isolated Morse sectors.

\subsection{The compactifiable \texorpdfstring{$B_2$}{B2} pencil}

For the projective calculation we use the standard representative
\begin{equation}
\label{eq:B2_compact_torus}
f_{B_2}(x,y,z)=\frac{(x+y+1)^4}{xyz}+z.
\end{equation}
Its pencil is
\begin{equation}
\label{eq:B2_compact_pencil}
F_{B_2,\lambda}
=(X+Y+T)^4+XYZ^2-\lambda XYZT=0
\end{equation}
in $\mathbb P^3_{[X:Y:Z:T]}$ \cite{PrzyjalkowskiWeak,Przyjalkowski2026}.
Put $A=x+y+1$ on $U_T$.  The logarithmic critical equations on the dense
torus are
\[
\theta_xf_{B_2}=\frac{A^3(4x-A)}{xyz},
\qquad
\theta_yf_{B_2}=\frac{A^3(4y-A)}{xyz},
\qquad
\theta_zf_{B_2}=-\frac{A^4}{xyz}+z.
\]
At a critical point one has $A\neq0$, so the first two equations give
$4x=A=4y$.  Therefore $x=y=1/2$, and the last equation gives $z^2=64$.
The two critical points are thus
\[
q_\pm=(1/2,1/2,\pm8),
\qquad
f_{B_2}(q_\pm)=\pm16,
\]
or, projectively,
\[
q_\pm=[1:1:\pm16:2].
\]
Direct differentiation gives
\begin{equation}
\label{eq:B2_compact_hessian}
\det(\theta_i\theta_jf_{B_2})(q_\pm)=\pm512,
\end{equation}
so both points are ordinary double points of their fibers.

At $\lambda=0$, the reduced raw singular support can be read on the four
projective charts:
\begin{center}
\small
\begin{tblr}{
  width=0.91\textwidth,
  colspec={Q[c,wd=0.10\textwidth]Q[l,wd=0.43\textwidth]X[l]},
  row{1}={font=\bfseries},
  hline{1,2,Z}={0.4pt},
  rowsep=2pt,
}
Chart & Central equation & Reduced singular support \\
$U_T$ & $(x+y+1)^4+xyz^2$ & $x+y+1=z=0$ \\
$U_X$ & $(1+y+t)^4+yz^2$ & $1+y+t=z=0$ \\
$U_Y$ & $(x+1+t)^4+xz^2$ & $x+1+t=z=0$ \\
$U_Z$ & $(x+y+t)^4+xy$ & $x=y=t=0$ \\
\end{tblr}
\end{center}
Globally this is the disjoint union
\begin{equation}
\label{eq:B2_raw_central_support}
L=V(X+Y+T,Z)
\qquad\text{and}\qquad
p_\infty=[0:0:1:0].
\end{equation}
On the open part of $L$ where $XY\neq0$, the transverse equation is
formally $u^4+v^2$, hence of type $A_3$.  At $p_\infty$, the $U_Z$ equation
is $(x+y+t)^4+xy$; eliminating the nondegenerate pair $(x,y)$ by Kn\"orrer
periodicity again leaves the one-variable $A_3$ potential $t^4$.  The two
special boundary points of $L$ and the base-locus intersections must be
retained in the blowup calculation; the generic $A_3$ description alone is
not a global decomposition.

\subsection{The compactifiable \texorpdfstring{$V_{10}$}{V10} pencil}

For $V_{10}$ we use
\begin{equation}
\label{eq:V10_compact_torus}
f_{V_{10}}(x,y,z)
=
\frac{(1+x+y+z+xy+xz+yz)^2}{xyz}.
\end{equation}
Write
\[
Q=T^2+T(X+Y+Z)+XY+XZ+YZ.
\]
The projective pencil is
\begin{equation}
\label{eq:V10_compact_pencil}
F_{V_{10},\lambda}=Q^2-\lambda XYZT=0
\end{equation}
\cite{PrzyjalkowskiWeak,Przyjalkowski2026}.  On $U_T$, let
\[
q=1+x+y+z+xy+xz+yz.
\]
The entire smooth surface $V(q)$ is critical in the central fiber.  Away
from it, the logarithmic critical equations are
\begin{equation}
\label{eq:V10_saturated_critical_equations}
2xq_x=q,
\qquad
2yq_y=q,
\qquad
2zq_z=q.
\end{equation}
After saturation by $xyzq$, their reduced solution is
\[
x=y=z=r,
\qquad
r^2-r-1=0.
\]
Writing $a=r^{-1}$ gives $a^2+a=1$.  The two noncentral critical points and
their values are therefore
\[
r_a=(a^{-1},a^{-1},a^{-1}),
\qquad
f_{V_{10}}(r_a)=20a+32,
\]
or, projectively,
\begin{equation}
\label{eq:V10_compact_nodes}
r_a=[1:1:1:a],
\qquad
a^2+a=1.
\end{equation}
Their logarithmic Hessian determinants reduce modulo $a^2+a-1$ to
\begin{equation}
\label{eq:V10_compact_hessian}
\det(\theta_i\theta_jf_{V_{10}})(r_a)
=32(123a+199)\neq0.
\end{equation}
Thus these are precisely two ordinary double points away from the base locus.

The matrix of the projective quadratic form $Q$ has determinant $1/16$.
Consequently the quadric
\[
\mathcal Q=V(Q)\subset\mathbb P^3
\]
is smooth, while the raw central member of
\eqref{eq:V10_compact_pencil} is the nonreduced divisor $2\mathcal Q$.
Away from $V(XYZT)$ its formal transverse equation is $u^2$, with the
coefficient twisting inherited from $1/(xyz)$.  The intersections with the
coordinate boundary are part of the base locus and cannot be discarded
before resolving the pencil.

\subsection{Resolved compactifications and the \texorpdfstring{$d=2$}{d=2} comparison}

Let
\[
\pi_F:Y_F\longrightarrow
\mathbb P^3\times\mathbb A^1_\lambda,
\qquad F\in\{B_2,V_{10}\},
\]
denote the standard crepant resolution of the corresponding pencil, and let
$w_F:Y_F\to\mathbb A^1$ be the resolved Landau--Ginzburg map.  Put
\[
\Gamma_F
=
\bigl(\CritLoc(w_F)\cap w_F^{-1}(0)\bigr)_{\mathrm{red}}.
\]
We define the compactified central categories by
\begin{equation}
\label{eq:d2_compact_central_categories}
\CcompBtwo
:=
\overline{\MF^{\Gamma_{B_2}}
(\widehat{Y_{B_2}}_{\Gamma_{B_2}},\widehat w_{B_2})},
\qquad
\CcompVten
:=
\overline{\MF^{\Gamma_{V_{10}}}
(\widehat{Y_{V_{10}}}_{\Gamma_{V_{10}}},\widehat w_{V_{10}})}.
\end{equation}
The standard resolution calculation gives
\begin{equation}
\label{eq:d2_component_count}
k(B_2)=h^{1,2}(B_2)+1=11,
\qquad
k(V_{10})=h^{1,2}(V_{10})+1=11
\end{equation}
for the numbers of irreducible components of the resolved central fibers
\cite{PrzyjalkowskiHodge}.  The noncentral audit above, together with the
base-locus resolution, gives no additional noncentral singularities beyond
the two ordinary double points on either side
\cite{Przyjalkowski2026}.  Hence
\begin{equation}
\label{eq:B2_compact_final_decomposition}
\overline{\MF(Y_{B_2},w_{B_2})}
\simeq
\CcompBtwo\oplus\MorseCat{3}^{\oplus2},
\end{equation}
and
\begin{equation}
\label{eq:V10_compact_final_decomposition}
\overline{\MF(Y_{V_{10}},w_{V_{10}})}
\simeq
\CcompVten\oplus\MorseCat{3}^{\oplus2}.
\end{equation}

\begin{proposition}[What the $d=2$ calculation proves]
\label{prop:d2_compactified_comparison}
For the standard compactifiable representatives of $B_2$ and $V_{10}$, the
isolated compactified sectors are equivalent:
\[
\mathcal I_{B_2}^{\mathrm{comp}}
\simeq
\mathcal I_{V_{10}}^{\mathrm{comp}}
\simeq
\MorseCat{3}^{\oplus2}.
\]
Their resolved central fibers also have the same number, eleven, of
irreducible components.  These equalities do not imply
\[
\CcompBtwo\simeq\CcompVten.
\]
Indeed, the raw central supports have different geometry: the $B_2$ support
is the line-plus-point configuration in
\eqref{eq:B2_raw_central_support}, whereas the $V_{10}$ central member is a
double smooth quadric.  An equivalence of the full central categories would
require a comparison of all exceptional components and their attaching
functors, and no such equivalence is obtained from the present calculation.
\end{proposition}

\begin{proof}
The first assertion follows from
\eqref{eq:B2_compact_hessian},
\eqref{eq:V10_compact_hessian},
\eqref{eq:B2_compact_final_decomposition}, and
\eqref{eq:V10_compact_final_decomposition}.  The component count is
\eqref{eq:d2_component_count}.  The final statement follows because neither
the number of components nor the number of isolated nodes determines the
formal gluing diagram of the resolved central fiber.
\end{proof}

\section{The \texorpdfstring{$d=3$}{d=3} pair: \texorpdfstring{$B_3$ and $V_{14}$}{B3 and V14}}
\label{sec:B3_V14_calculation}

We next treat the cubic threefold $B_3$ and the prime Fano threefold
$V_{14}$.  This pair illustrates a second boundary phenomenon.  The two
isolated contributions of $B_3$ are already visible on the dense torus.  For
the convenient $V_{14}$ compactification, one node is visible on the torus,
whereas the other appears only after resolving a curve in the base locus.

\subsection{The source Laurent polynomials on the torus}

The source polynomial for $B_3$ is
\begin{equation}
\label{eq:B3_source_polynomial}
W_{B_3}
=x+\frac{x}{yz}+y+z+\frac2z+\frac2y
+\frac{y}{xz}+\frac2x+\frac{z}{xy}.
\end{equation}
Its saturated logarithmic critical ideal is
\[
(x-y,\ y-z,\ z^2-3).
\]
Thus the two critical points are $(a,a,a)$, where $a^2=3$, and
\[
W_{B_3}(a,a,a)=6a,
\qquad
\det(\theta_i\theta_jW_{B_3})(a,a,a)=54a\neq0.
\]
Consequently,
\begin{equation}
\label{eq:B3_source_category}
\overline{\MF((\mathbb C^*)^3,W_{B_3})}
\simeq
\bigoplus_{a^2=3}\MorseCat{3}.
\end{equation}

For the source $V_{14}$ polynomial, put
\[
A=yz^2+yz+1
\]
and
\[
\begin{aligned}
B={}&x^2yz^2+xy^2z^4+2xy^2z^3+xy^2z^2
     +2xyz+x+yz^2.
\end{aligned}
\]
The exact factorization
\begin{equation}
\label{eq:V14_source_factorization}
W_{V_{14}}+4=\frac{AB}{xy^2z^3}
\end{equation}
shows that the critical fiber at $-4$ contains the smooth curve
\[
\Gamma_{\mathrm{src}}
=
\left\{
\left(1,-\frac1{t(t+1)},t\right):
t\in\mathbb C^*\setminus\{-1\}
\right\}.
\]
Indeed, on $A=0$ one has
\[
y=-\frac1{z(z+1)},
\qquad
B=-\frac{z(x-1)^2}{z+1}.
\]
With $R=\mathbb C[t^{\pm1},(t+1)^{-1}]$ and global normal coordinates
\[
U=A,
\qquad
V=x-1,
\]
the completed shifted potential is, after multiplication by a unit and a
relative finite-determinacy change of coordinates,
\[
U\left(U^2-\frac{t}{t+1}V^2\right).
\]
The coefficient has square class
\[
\left[\frac{t}{t+1}\right]=[t(t+1)]
\quad\text{in }R^*/R^{*2}.
\]
The two complementary critical points are
\[
P_+=(1,1,1),
\qquad
P_-=\left(-1,9,-\frac13\right),
\]
with
\[
W_{V_{14}}(P_+)=23,
\qquad
W_{V_{14}}(P_-)=-5,
\]
and logarithmic Hessian determinants $1458$ and $-2$, respectively.  Hence
\begin{equation}
\label{eq:V14_source_category}
\overline{\MF((\mathbb C^*)^3,W_{V_{14}})}
\simeq
\Dper{\Gamma_{\mathrm{src}}}{4}{t(t+1)}
\oplus\MorseCat{3}^{\oplus2}.
\end{equation}
In particular, the source categories in
\eqref{eq:B3_source_category} and \eqref{eq:V14_source_category} are not
equivalent: the second contains a positive-dimensional residual $D_4$
contribution absent from the first.

\subsection{The compactifiable \texorpdfstring{$B_3$}{B3} pencil}

For the compactification calculation, take
\begin{equation}
\label{eq:B3_compact_torus}
f_{B_3}(x,y,z)=\frac{(x+y+1)^3}{xyz}+z.
\end{equation}
Writing $A_h=X+Y+T$, the associated quartic pencil is
\begin{equation}
\label{eq:B3_compact_pencil}
F_{B_3,\lambda}
=T A_h^3+XYZ^2-\lambda XYZT=0
\end{equation}
\cite{PrzyjalkowskiWeak,Przyjalkowski2026}.  On the dense torus, set
$A=x+y+1$.  The logarithmic derivatives are
\[
\theta_xf_{B_3}=\frac{A^2(3x-A)}{xyz},
\qquad
\theta_yf_{B_3}=\frac{A^2(3y-A)}{xyz},
\qquad
\theta_zf_{B_3}=-\frac{A^3}{xyz}+z.
\]
At a critical point $A\neq0$, so $3x=A=3y$.  It follows that $x=y=1$ and
$z^2=27$.  If $a^2=3$, the two points, their values, and their Hessians are
\begin{equation}
\label{eq:B3_compact_nodes}
q_a=(1,1,3a),
\qquad
f_{B_3}(q_a)=6a,
\qquad
\det(\theta_i\theta_jf_{B_3})(q_a)=54a.
\end{equation}
Thus both are ordinary double points.  Projectively they are
\[
q_a=[1:1:3a:1].
\]

At $\lambda=0$, the four affine equations and their reduced singular
supports are
\begin{center}
\small
\begin{tblr}{
  width=0.91\textwidth,
  colspec={Q[c,wd=0.10\textwidth]Q[l,wd=0.43\textwidth]X[l]},
  row{1}={font=\bfseries},
  hline{1,2,Z}={0.4pt},
  rowsep=2pt,
}
Chart & Central equation & Reduced singular support \\
$U_T$ & $(x+y+1)^3+xyz^2$ & $x+y+1=z=0$ \\
$U_X$ & $t(1+y+t)^3+yz^2$ & $1+y+t=z=0$ \\
$U_Y$ & $t(x+1+t)^3+xz^2$ & $x+1+t=z=0$ \\
$U_Z$ & $t(x+y+t)^3+xy$ & $x=y=t=0$ \\
\end{tblr}
\end{center}
Thus the raw central support is again
\begin{equation}
\label{eq:B3_raw_central_support}
L_3=V(X+Y+T,Z)
\qquad\text{together with}\qquad
p_\infty=[0:0:1:0].
\end{equation}
On the open part of $L_3$ where the coefficients $T$ and $XY$ are units, the
transverse equation is $u^3+v^2$, of type $A_2$.  At $p_\infty$, the equation
$t(x+y+t)^3+xy$ has a nondegenerate split quadratic pair $(x,y)$; eliminating
it leaves a one-variable potential with leading term $t^4$, hence type $A_3$.
As before, the special points of $L_3$ and its base-locus intersections must
be handled by the global resolution.

\subsection{The compactifiable \texorpdfstring{$V_{14}$}{V14} pencil}

For $V_{14}$ use the standard representative
\begin{equation}
\label{eq:V14_compact_torus}
f_{V_{14}}
=
\frac{(x+y+z+1)^2}{x}
+
\frac{(x+y+z+1)(y+z+1)(z+1)^2}{xyz}.
\end{equation}
Set
\[
A_h=X+Y+Z+T,
\qquad
B_h=Y+Z+T,
\qquad
C_h=Z+T.
\]
The projective pencil is
\begin{equation}
\label{eq:V14_compact_pencil}
F_{V_{14},\lambda}
=A_h\bigl(YZ A_h+B_hC_h^2\bigr)-\lambda XYZT=0
\end{equation}
\cite{PrzyjalkowskiWeak,Przyjalkowski2026}.  Its dense-torus critical locus
contains the central curve
\begin{equation}
\label{eq:V14_compact_central_curve}
\Gamma_{\mathrm{cmp}}
=
\{(x,y,z)=(t,-t,-1):t\in\mathbb C^*\}.
\end{equation}
There is one complementary torus critical point
\[
r=(6,3/2,1/2),
\qquad
r=[12:3:1:2]
\]
with
\begin{equation}
\label{eq:V14_visible_node}
f_{V_{14}}(r)=27,
\qquad
\det(\theta_i\theta_jf_{V_{14}})(r)=1458.
\end{equation}
This gives the first noncentral ordinary double point.

The raw central member factors as
\begin{equation}
\label{eq:V14_raw_central_factorization}
F_{V_{14},0}=A_hG_h,
\qquad
G_h=YZ A_h+B_hC_h^2.
\end{equation}
Thus it is the union of the plane $\Pi=V(A_h)$ and the cubic
$D=V(G_h)$.  Their scheme-theoretic intersection on $\Pi$ is
\[
A_h=B_hC_h^2=0;
\]
the component $V(A_h,C_h)$ occurs with multiplicity two and restricts to
$\Gamma_{\mathrm{cmp}}$ on the dense torus.  To identify the transverse
singularity, use
\[
t=x,
\qquad
U=x+y+z+1,
\qquad
V=z+1,
\qquad
y=U-t-V.
\]
Then the potential is exactly
\begin{equation}
\label{eq:V14_compact_local_equation}
f_{V_{14}}
=
\frac{U\bigl(-U^2V+U^2+(t-1)UV-tU+tV^2\bigr)}
{t(V-1)(-U+V+t)}.
\end{equation}
Since the denominator is a unit along $U=V=0$, the transverse germ is
contact-equivalent to
\[
U(U-V^2)
=
\left(U-\frac12V^2\right)^2-\frac14V^4,
\]
an $A_3$ family along $\Gamma_{\mathrm{cmp}}$.  This differs from the
residual $D_4$ family of the source representative
\eqref{eq:V14_source_factorization}.

\subsection{The boundary node of \texorpdfstring{$V_{14}$}{V14}}

The second isolated contribution is created by resolving the base curve
through
\[
P_6=[0:1:-1:0].
\]
On $U_Y$, write $x=X/Y$, $z=Z/Y$, $t=T/Y$, and set
\[
b=1+z+t.
\]
The relevant base curve is $x=b=0$.  In the blowup chart $x=br$, after
dividing by the exceptional factor $b$, the strict transform of the pencil is
\begin{equation}
\label{eq:V14_first_blowup_equation}
\widetilde F_\lambda
=
b(r+1)^2(b-t-1)
+b(r+1)(b-1)^2
-\lambda r(b-t-1)t.
\end{equation}
For every $\lambda\neq0$, the point $b=r=t=0$ is singular.  Its Hessian is
\[
H_\lambda=
\begin{pmatrix}
-2&-1&-1\\
-1&0&\lambda\\
-1&\lambda&0
\end{pmatrix},
\qquad
\det H_\lambda=2\lambda(\lambda+1).
\]
Hence this is an ordinary double point for $\lambda\neq0,-1$, while the
fiber at $\lambda=-1$ has the more degenerate $A_3$ point that survives the
simultaneous resolution step.

To see the remaining node explicitly, set $\lambda=-1$ and blow up
$b=r=t=0$.  In the chart
\[
b=e,
\qquad
r=eR,
\qquad
t=eT,
\]
the strict transform is
\begin{equation}
\label{eq:V14_second_blowup_equation}
\begin{aligned}
G(e,R,T)={}&-R^2Te^2+R^2e^2-R^2e-RT^2e-RTe-RT\\
&+Re^2-R-T+e-1.
\end{aligned}
\end{equation}
It has a critical point at
\[
(e,R,T)=(0,-1,-1),
\]
where
\[
\operatorname{Hess}G=
\begin{pmatrix}
2&2&-1\\
2&0&-1\\
-1&-1&0
\end{pmatrix},
\qquad
\det(\operatorname{Hess}G)=2.
\]
This is the second ordinary double point, lying in the fiber at
$\lambda=-1$.  The calculation also explains why it is absent from the
original dense-torus audit \cite{Przyjalkowski2026}.

\subsection{Resolved compactifications and the \texorpdfstring{$d=3$}{d=3} comparison}

Let $w_F:Y_F\to\mathbb A^1$ be the standard resolved compactification for
$F\in\{B_3,V_{14}\}$, and put
\[
\Gamma_F^{(3)}
=
\bigl(\CritLoc(w_F)\cap w_F^{-1}(0)\bigr)_{\mathrm{red}}.
\]
Define
\begin{equation}
\label{eq:d3_compact_central_categories}
\CcompBthree
:=
\overline{\MF^{\Gamma_{B_3}^{(3)}}
(\widehat{Y_{B_3}}_{\Gamma_{B_3}^{(3)}},\widehat w_{B_3})},
\qquad
\CcompVfourteen
:=
\overline{\MF^{\Gamma_{V_{14}}^{(3)}}
(\widehat{Y_{V_{14}}}_{\Gamma_{V_{14}}^{(3)}},\widehat w_{V_{14}})}.
\end{equation}
The central component formula gives
\begin{equation}
\label{eq:d3_component_count}
k(B_3)=h^{1,2}(B_3)+1=6,
\qquad
k(V_{14})=h^{1,2}(V_{14})+1=6
\end{equation}
\cite{PrzyjalkowskiHodge}.  Combining the two torus nodes
\eqref{eq:B3_compact_nodes} with the visible and boundary nodes
\eqref{eq:V14_visible_node}--\eqref{eq:V14_second_blowup_equation} gives
\begin{equation}
\label{eq:B3_compact_final_decomposition}
\overline{\MF(Y_{B_3},w_{B_3})}
\simeq
\CcompBthree\oplus\MorseCat{3}^{\oplus2},
\end{equation}
and
\begin{equation}
\label{eq:V14_compact_final_decomposition}
\overline{\MF(Y_{V_{14}},w_{V_{14}})}
\simeq
\CcompVfourteen\oplus\MorseCat{3}^{\oplus2}.
\end{equation}

\begin{proposition}[What the $d=3$ calculation proves]
\label{prop:d3_compactified_comparison}
For the standard compactifications of $B_3$ and $V_{14}$,
\[
\mathcal I_{B_3}^{\mathrm{comp}}
\simeq
\mathcal I_{V_{14}}^{\mathrm{comp}}
\simeq
\MorseCat{3}^{\oplus2},
\]
and both resolved central fibers have six irreducible components.  On the
$V_{14}$ side, one Morse summand is supported at the torus point of value
$27$, and the other is the boundary node of value $-1$ produced by the
base-curve resolution.  These facts do not by themselves prove
\[
\CcompBthree\simeq\CcompVfourteen.
\]
Such an equivalence requires compatible global attaching data, or the
compactified-HMS identifications for a Kuznetsov-matched pair.
\end{proposition}

\begin{proof}
The isolated statement follows from
\eqref{eq:B3_compact_nodes}, \eqref{eq:V14_visible_node}, and the explicit
boundary calculation
\eqref{eq:V14_first_blowup_equation}--\eqref{eq:V14_second_blowup_equation}.
The component count is \eqref{eq:d3_component_count}.  Equality of these
numerical and isolated sectors does not identify the formal gluing diagrams
of the six central components.
\end{proof}

\section{The remaining index-two compactifications:
\texorpdfstring{$B_1$ and $B_5$}{B1 and B5}}
\label{sec:B1_B5_compactifications}

We now complete the compactification calculation on the index-two side.  In
the notation of \cite{Przyjalkowski2026}, our $B_d$ is the threefold denoted
$Y_d$.  The cases $B_2$, $B_3$, and $B_4$ were calculated above; here we work
directly with the projective pencils for $B_1$ and $B_5$, without introducing
or solving separate source-torus models.

\subsection{The compactified \texorpdfstring{$B_1$}{B1} pencil}

Put
\[
H=XY-XT-T^2.
\]
A standard toric Landau--Ginzburg model for $B_1$ is birational over the base
to the quartic pencil
\begin{equation}
\label{eq:B1_compact_pencil}
F_{B_1,\lambda}
=X^4+Z^2H-\lambda YZH=0
\end{equation}
in $\mathbb P^3_{[X:Y:Z:T]}$
\cite{PrzyjalkowskiWeak,Przyjalkowski2026}.  The singular points belonging
to the base-locus resolution are
\[
\begin{gathered}
P_1=[0:0:0:1],\qquad P_2=[0:1:0:0],\\
P_3=[0:0:1:0],\qquad P_4=[0:1:\lambda:0].
\end{gathered}
\]
For $\lambda\neq0$, their surface singularity types are, respectively,
$A_3$, $A_7$, $A_1$, and $A_7$.  In the chart $Y=1$, set
$X'=x(1-t)-t^2$ and translate the local $z$-coordinate to $Z'$.  The
equation becomes
$Z'X'+((X'+t^2)/(1-t))^4$; its restriction to the Hessian kernel has
residual order eight.  The same order occurs at $P_4$.  They are resolved as part of the standard
base-locus resolution and do not contribute isolated singularities to the
resolved noncentral fibers.

The two singularities away from the base locus can be obtained directly on
$U_T$.  Write
\[
h=xy-x-1,
\qquad
f_\lambda=x^4+(z^2-\lambda yz)h.
\]
At a non-base singular point, $z\neq0$ and $h\neq0$.  The equations
$\partial_zf_\lambda=\partial_yf_\lambda=0$ give
\[
z=\frac{\lambda y}{2},
\qquad
zx+\lambda h=0.
\]
Combining these identities with $f_\lambda=\partial_xf_\lambda=0$ yields
\[
x=-2,\qquad y=\frac13,\qquad z=4\alpha,\qquad
\lambda=24\alpha,\qquad \alpha^2=3.
\]
Thus the two projective points are
\begin{equation}
\label{eq:B1_compact_nodes}
q_\alpha=[-6:1:12\alpha:3],
\qquad
\lambda_\alpha=24\alpha,
\qquad
\alpha^2=3.
\end{equation}
The determinant of the affine Hessian in the variables $(x,y,z)$ is
\begin{equation}
\label{eq:B1_compact_hessian}
\det\bigl(\operatorname{Hess}_{x,y,z}f_{\lambda_\alpha}
(q_\alpha)\bigr)=13824,
\end{equation}
so both points are ordinary double points.

For completeness, the reduced singular support of the raw central member is
\begin{equation}
\label{eq:B1_raw_central_support}
\Sing(F_{B_1,0})_{\mathrm{red}}
=V(X,Z)\ \sqcup\ \{[0:0:1:0]\}.
\end{equation}
At a general point of the line $V(X,Z)$ the transverse equation is
$u^4+v^2$, of type $A_3$.  The isolated point lies on the base locus and is
removed by its crepant resolution; it is not a third isolated summand of the
resolved Landau--Ginzburg model.

Let $w_{B_1}:Y_{B_1}\to\mathbb A^1$ be the standard resolved
compactification and put
\[
\Gamma_{B_1}
=\bigl(\CritLoc(w_{B_1})\cap w_{B_1}^{-1}(0)\bigr)_{\mathrm{red}}.
\]
Define
\begin{equation}
\label{eq:B1_compact_central_category}
\CcompBone
:=
\overline{\MF^{\Gamma_{B_1}}
(\widehat{Y_{B_1}}_{\Gamma_{B_1}},\widehat w_{B_1})}.
\end{equation}
The resolved central fiber has
\begin{equation}
\label{eq:B1_component_count}
k(B_1)=h^{1,2}(B_1)+1=22
\end{equation}
irreducible components \cite{PrzyjalkowskiHodge}.  Equations
\eqref{eq:B1_compact_nodes} and \eqref{eq:B1_compact_hessian}, together with
the base-locus resolution, give
\begin{equation}
\label{eq:B1_compact_final_decomposition}
\overline{\MF(Y_{B_1},w_{B_1})}
\simeq
\CcompBone\oplus\MorseCat{3}^{\oplus2}.
\end{equation}

\begin{remark}[Two corrected substitutions]
\label{rem:corrected_B1_B5_data}
Direct substitution into the displayed pencils corrects two typographical
errors in the point data printed in \cite{Przyjalkowski2026}.  For $B_1$ the
node has first homogeneous coordinate $-6$, as in
\eqref{eq:B1_compact_nodes}.  For $B_5$ below, the critical value is
$4\alpha^3+6\alpha$.  The corrected expressions make the defining equation
and all first derivatives vanish; the Hessian calculations then verify the
ordinary-double-point assertions.
\end{remark}

\subsection{The compactified \texorpdfstring{$B_5$}{B5} pencil}

For $B_5$, the standard quartic pencil is
\begin{equation}
\label{eq:B5_compact_pencil}
\begin{aligned}
F_{B_5,\lambda}={}&
X^2YZ+XY^2Z+XYZ^2+YZT^2+XZT^2+XYT^2+T^4\\
&-\lambda XYZT=0.
\end{aligned}
\end{equation}
Its base-locus singular points are
\[
P_1^1=[1:0:0:0],\quad
P_1^2=[0:1:0:0],\quad
P_1^3=[0:0:1:0]
\]
of type $A_4$ for $\lambda\neq0$ and type $A_5$ on the central member
$\lambda=0$, whereas the points
\[
P_2^1=[1:-1:0:0],\quad
P_2^2=[1:0:-1:0],\quad
P_2^3=[0:1:-1:0]
\]
are of type $A_1$.  In the chart $X=1$, a local equation at $P_1^1$ is
$(y+t^2)(z+t^2)+yz(y+z+t^2-\lambda t)$.  After eliminating the
nondegenerate quadratic directions, the residual term begins with
$-\lambda t^5$.  Thus it has order five for $\lambda\neq0$ and order six
for $\lambda=0$; symmetry gives the same types at the other $P_1^i$.
Resolving these fixed singularities resolves the base
locus without creating any further singular fibers
\cite{Przyjalkowski2026}.

The remaining singularities are found in the chart $U_X$.  Symmetry and the
three first-derivative equations give
\begin{equation}
\label{eq:B5_compact_nodes}
p_\alpha=[1:1:1:\alpha],
\qquad
\alpha^4+\alpha^2=1,
\qquad
\lambda_\alpha=4\alpha^3+6\alpha.
\end{equation}
There are four roots of $\alpha^4+\alpha^2-1$, hence four such points.  A
direct calculation in the affine variables $(Y/X,Z/X,T/X)$ gives
\begin{equation}
\label{eq:B5_compact_hessian}
\det\bigl(\operatorname{Hess}F_{B_5,\lambda_\alpha}
(p_\alpha)\bigr)
=2(3\alpha^2+4)\neq0.
\end{equation}
Thus all four points are ordinary double points away from the base locus.
The resolved pencil is Lefschetz: it has four singular fibers, each with one
ordinary double point \cite{Przyjalkowski2026}.  Since
$h^{1,2}(B_5)=0$, the fiber over $0$ is smooth and irreducible after the
base-locus resolution, so
\begin{equation}
\label{eq:B5_component_count}
k(B_5)=1,
\qquad
\mathcal C^{\mathrm{comp}}_{B_5,0}=0.
\end{equation}
Consequently,
\begin{equation}
\label{eq:B5_compact_final_decomposition}
\overline{\MF(Y_{B_5},w_{B_5})}
\simeq
\MorseCat{3}^{\oplus4}.
\end{equation}

\begin{proposition}[Compactification census for all five $B_d$]
\label{prop:Bd_compactification_census}
For the standard resolved compactifications of the index-two Fano
threefolds $B_d$, $1\leq d\leq5$, one has
\[
\begin{array}{c|c|c|c}
d & \varepsilon(B_d) & k(B_d) &
\overline{\MF(Y_{B_d},w_{B_d})} \\
\hline
1 & 2 & 22 & \CcompBone\oplus\MorseCat{3}^{\oplus2} \\
2 & 2 & 11 & \CcompBtwo\oplus\MorseCat{3}^{\oplus2} \\
3 & 2 & 6  & \CcompBthree\oplus\MorseCat{3}^{\oplus2} \\
4 & 2 & 3  & \CcompB\oplus\MorseCat{3}^{\oplus2} \\
5 & 4 & 1  & \MorseCat{3}^{\oplus4}.
\end{array}
\]
Here $\varepsilon(B_d)$ is the number of compactified ordinary double
points, and $k(B_d)$ is the number of irreducible components of the resolved
fiber over $0$.  For $d\leq4$, the notation $\mathcal C^{\mathrm{comp}}_{B_d,0}$
denotes the connected formal category supported on that central fiber.  The
table is a compactification statement and does not require a separate
calculation of the critical locus of a source Laurent polynomial on the
torus.
\end{proposition}

\begin{proof}
The rows $d=2,3,4$ are
\eqref{eq:B2_compact_final_decomposition},
\eqref{eq:B3_compact_final_decomposition}, and
\eqref{eq:B4_final_decomposition}.  The rows $d=1,5$ are
\eqref{eq:B1_compact_final_decomposition} and
\eqref{eq:B5_compact_final_decomposition}.  The component counts follow from
$k(B_d)=h^{1,2}(B_d)+1$ for $d\leq4$ and from the smooth irreducible central
fiber for the Lefschetz pencil $B_5$.
\end{proof}

\section{The remaining prime compactifications:
\texorpdfstring{$V_6$ and $V_{22}$}{V6 and V22}}
\label{sec:V6_V22_compactifications}

We now calculate the two extreme cases that were still missing on the
index-one side.  As for $B_1$ and $B_5$, we start directly from the
compactified pencils.  Thus no separate census of the critical locus of a
source Laurent polynomial is needed.

\subsection{The compactified \texorpdfstring{$V_6$}{V6} pencil}

The threefold $V_6$ is a complete intersection of a quadric and a cubic.
Its Givental-type model has the anticanonical compactification
\begin{equation}
\label{eq:V6_compact_pencil}
(X_0+X_1)^2(Y+Z+T)^3
=\lambda X_0X_1YZT
\end{equation}
in
$\mathbb P^1_{[X_0:X_1]}\times\mathbb P^2_{[Y:Z:T]}$.
The equation has bidegree $(2,3)$, which is the anticanonical degree of the
ambient product \cite{PrzyjalkowskiWeak,PrzyjalkowskiCY}.

On the chart $X_0=T=1$, put $x=X_1/X_0$, $y=Y/T$, and $z=Z/T$.  Away from
the base locus, the logarithmic derivative equations for the quotient of the
two pencil generators are
\[
\frac{2x}{x+1}=1,
\qquad
\frac{3y}{y+z+1}=1,
\qquad
\frac{3z}{y+z+1}=1.
\]
They have the unique solution
\begin{equation}
\label{eq:V6_compact_node}
(x,y,z)=(1,1,1),
\qquad
\lambda=108.
\end{equation}
For
\[
f_{V_6,\lambda}
=(x+1)^2(y+z+1)^3-\lambda xyz,
\]
the Hessian at this point is
\begin{equation}
\label{eq:V6_compact_hessian}
\operatorname{Hess}f_{V_6,108}(1,1,1)
=
\begin{pmatrix}
54&0&0\\
0&72&-36\\
0&-36&72
\end{pmatrix},
\qquad
\det=209952=2^5 3^8.
\end{equation}
Hence the noncentral singularity is an ordinary double point.  The standard
crepant resolution of the fixed base locus creates no further noncentral
singularities; equivalently, the complete-intersection calculation gives
exactly one such fiber \cite{Przyjalkowski2026}.

The raw member over $0$ is the nonreduced divisor
\begin{equation}
\label{eq:V6_raw_central_member}
2D_X+3D_A,
\qquad
D_X=V(X_0+X_1),\qquad D_A=V(Y+Z+T).
\end{equation}
Let $w_{V_6}:Y_{V_6}\to\mathbb A^1$ be the resolved compactification and
write
\[
\Gamma_{V_6}
=\bigl(\CritLoc(w_{V_6})\cap w_{V_6}^{-1}(0)\bigr)_{\mathrm{red}}.
\]
We set
\begin{equation}
\label{eq:V6_compact_central_category}
\CcompVsix
:=
\overline{\MF^{\Gamma_{V_6}}
(\widehat{Y_{V_6}}_{\Gamma_{V_6}},\widehat w_{V_6})}.
\end{equation}
Since $h^{1,2}(V_6)=20$, the resolved central fiber has
\begin{equation}
\label{eq:V6_component_count}
k(V_6)=21
\end{equation}
irreducible components.  The compactified decomposition is therefore
\begin{equation}
\label{eq:V6_compact_final_decomposition}
\overline{\MF(Y_{V_6},w_{V_6})}
\simeq
\CcompVsix\oplus\MorseCat{3}.
\end{equation}

\subsection{The compactified \texorpdfstring{$V_{22}$}{V22} pencil}

For $V_{22}$, homogenizing the standard Laurent representative, including
its constant term $4$ and the monomial $z/(xy)$, gives the quartic pencil
\begin{equation}
\label{eq:V22_compact_pencil}
\begin{aligned}
F_{V_{22},\lambda}={}&
X^2Y^2+XY^2T+X^2YT+X^2YZ+XY^2Z+XYT^2+4XYZT\\
&+YZT^2+XZT^2+XYZ^2+ZT^3+YZ^2T+XZ^2T+Z^2T^2\\
&-\lambda XYZT=0
\end{aligned}
\end{equation}
in $\mathbb P^3_{[X:Y:Z:T]}$
\cite{PrzyjalkowskiWeak,Przyjalkowski2026}.  Let $\alpha$ satisfy
\[
\alpha^3+2\alpha^2-2=0.
\]
Direct substitution gives three singular points away from the base locus,
\begin{equation}
\label{eq:V22_visible_nodes}
p_\alpha=[1:1:1:\alpha],
\qquad
\lambda_\alpha=3\alpha^2+8\alpha+8.
\end{equation}
The cubic has three distinct roots.  On the chart $X=1$, the Hessian
determinant at $p_\alpha$ reduces modulo
$\alpha^3+2\alpha^2-2$ to
\begin{equation}
\label{eq:V22_visible_hessian}
2(14\alpha^2+39\alpha+32).
\end{equation}
Moreover,
\[
\operatorname{Res}_\alpha
(\alpha^3+2\alpha^2-2,14\alpha^2+39\alpha+32)=22,
\]
so all three points are ordinary double points.

The singular points that remain on the fixed base locus are
\[
\begin{gathered}
P_1^1=[1:0:0:0],\quad P_1^2=[0:1:0:0],\\
P_2^1=[1:0:0:-1],\quad P_2^2=[0:1:0:-1],\\
P_3=[0:0:1:0],\quad P_4=[1:1:-1:0],\quad
P_5=[1:1:0:-1],\\
P_5^1=[1:0:1:-1],\qquad P_5^2=[0:1:1:-1].
\end{gathered}
\]
The points other than $P_4$ are removed by the standard resolution of the
fixed base components.  The point $P_4$, however, produces one additional
ordinary double point.  We make this boundary contribution explicit.

Work on $X=1$ and set
\[
q=Y/X,\qquad u=(Y+Z)/X,\qquad t=T/X.
\]
The base line through $P_4$ is $u=t=0$.  In the blow-up chart $u=tr$, the
total transform of \eqref{eq:V22_compact_pencil} has a common factor $t$.
After removing it, the strict transform is
\begin{equation}
\label{eq:V22_boundary_strict_transform}
\begin{aligned}
g_\lambda={}&
\lambda q^2+q^3-q^2r-2q^2+qr+q
+qrt(-\lambda-2q+r+2)\\
&+t^2(qr^2-qr-q+r^2+r)+r(r+1)t^3.
\end{aligned}
\end{equation}
At $\lambda=0$, the point
\begin{equation}
\label{eq:V22_boundary_node}
(q,r,t)=(1,0,0)
\end{equation}
is singular.  The transformed second pencil generator is
$q(tr-q)$, whose value there is $-1$; hence the point is no longer on the
base locus.  Finally,
\begin{equation}
\label{eq:V22_boundary_hessian}
\operatorname{Hess}g_0(1,0,0)
=
\begin{pmatrix}
2&-1&0\\
-1&0&0\\
0&0&-2
\end{pmatrix},
\qquad
\det=2.
\end{equation}
Thus the blow-up creates precisely the fourth ordinary double point.  After
resolving the remaining fixed base singularities, the $V_{22}$ model is a
Lefschetz pencil with four singular fibers, one node in each
\cite{Przyjalkowski2026}.  Since $h^{1,2}(V_{22})=0$, there is no
positive-dimensional central contribution, and
\begin{equation}
\label{eq:V22_compact_final_decomposition}
k(V_{22})=1,
\qquad
\overline{\MF(Y_{V_{22}},w_{V_{22}})}
\simeq
\MorseCat{3}^{\oplus4}.
\end{equation}

\begin{remark}[The corrected $V_{22}$ homogenization]
\label{rem:corrected_V22_homogenization}
The transcription of the $V_{22}$ quartic in
\cite{Przyjalkowski2026} omits the terms $4XYZT$ and $Z^2T^2$ that arise
from the constant $4$ and $z/(xy)$ in the standard Laurent polynomial of
\cite{PrzyjalkowskiWeak}.  The constant term can be absorbed by translating
the parameter, but $Z^2T^2$ cannot.  Equation
\eqref{eq:V22_compact_pencil} is the direct homogenization; with these terms,
the points \eqref{eq:V22_visible_nodes}, all first derivatives, and the
boundary blow-up calculation vanish exactly as stated.
\end{remark}

\begin{proposition}[The two extreme prime compactifications]
\label{prop:V6_V22_compactification_census}
For the standard resolved compactifications one has
\[
\begin{array}{c|c|c|c}
F & \varepsilon(F) & k(F) & \overline{\MF(Y_F,w_F)}\\
\hline
V_6 & 1 & 21 & \CcompVsix\oplus\MorseCat{3}\\
V_{22} & 4 & 1 & \MorseCat{3}^{\oplus4}.
\end{array}
\]
For $V_6$, the nonisolated category is supported on the resolved central
fiber.  For $V_{22}$, all four contributions are isolated, including the
boundary node \eqref{eq:V22_boundary_node}.
\end{proposition}

\begin{proof}
The $V_6$ row is
\eqref{eq:V6_compact_hessian}--\eqref{eq:V6_compact_final_decomposition}.
The three non-base $V_{22}$ nodes are verified by
\eqref{eq:V22_visible_nodes}--\eqref{eq:V22_visible_hessian}; the fourth is
\eqref{eq:V22_boundary_node}--\eqref{eq:V22_boundary_hessian}.  The component
counts follow from $k(F)=h^{1,2}(F)+1$ when $h^{1,2}(F)>0$, and from the
irreducibility of the $V_{22}$ Lefschetz fibers.
\end{proof}

\section{Compactification comparison for all five pairs}
\label{sec:five_pair_compactification}

We now extend the compactification comparison from the detailed pair
$(B_4,V_{18})$ to
\[
(B_d,V_{4d+2}),
\qquad d=1,2,3,4,5.
\]
On the $B$ side, Proposition~\ref{prop:Bd_compactification_census} gives the
compactified calculation for every $d$.  Proposition
\ref{prop:V6_V22_compactification_census} supplies the two extreme cases on
the $V$ side.  For $d=2,3$,
Sections~\ref{sec:B2_V10_calculation} and
\ref{sec:B3_V14_calculation} additionally determine the raw central supports
and the two isolated nodes on the corresponding $V$ side.  For $d\neq4$,
however, we do not claim a complete chartwise computation of all exceptional
divisors and their attaching maps.

\subsection{Isolated compactification contributions}

Let
\[
\mu_{\mathrm{tor}}^{\mathrm{iso}}(F)
=\#\{\text{isolated Morse points of }W_F\text{ in }(\mathbb C^*)^3\},
\]
let $\varepsilon(F)$ be the number of ordinary double points in the chosen
compactification, and put
\[
\delta_\partial(F)
=\varepsilon(F)-\mu_{\mathrm{tor}}^{\mathrm{iso}}(F).
\]
Every isolated compactified singularity in this census is an ordinary double
point.  Consequently its matrix-factorization contribution is
\begin{equation}
\label{eq:five_pair_isolated_category}
\mathcal I_F^{\mathrm{comp}}
\simeq
\MorseCat{3}^{\oplus\varepsilon(F)}.
\end{equation}
The open-torus and compactified counts are as follows
\cite{Przyjalkowski2026}.  In this table, $\mu_B$ and $\mu_V$ count the
isolated Morse points on the dense torus of the compactifiable representatives
whose projective pencils are displayed; a different source or mutation
representative can have a different torus count.
\begin{center}
\footnotesize
\begin{tblr}{
  width=0.98\textwidth,
  colspec={Q[c,wd=0.05\textwidth]Q[c,wd=0.13\textwidth]
    Q[c,wd=0.08\textwidth]Q[c,wd=0.08\textwidth]
    Q[c,wd=0.08\textwidth]Q[c,wd=0.08\textwidth]
    Q[c,wd=0.09\textwidth]X[l]},
  colsep=2pt,
  row{1}={font=\bfseries},
  rowsep=2pt,
  hline{1,2,Z}={0.4pt},
}
$d$ & Pair & $\mu_B$ & $\varepsilon_B$ & $\mu_V$ & $\varepsilon_V$ &
$\delta_{\partial,V}$ & Isolated compactified categories \\
$5$ & $(B_5,V_{22})$ & $4$ & $4$ & $3$ & $4$ & $1$ &
$\mathcal I_{B_5}^{\mathrm{comp}}\simeq
 \mathcal I_{V_{22}}^{\mathrm{comp}}\simeq\MorseCat{3}^{\oplus4}$. \\
$4$ & $(B_4,V_{18})$ & $2$ & $2$ & $2$ & $2$ & $0$ &
$\mathcal I_{B_4}^{\mathrm{comp}}\simeq
 \mathcal I_{V_{18}}^{\mathrm{comp}}\simeq\MorseCat{3}^{\oplus2}$. \\
$3$ & $(B_3,V_{14})$ & $2$ & $2$ & $1$ & $2$ & $1$ &
$\mathcal I_{B_3}^{\mathrm{comp}}\simeq
 \mathcal I_{V_{14}}^{\mathrm{comp}}\simeq\MorseCat{3}^{\oplus2}$. \\
$2$ & $(B_2,V_{10})$ & $2$ & $2$ & $2$ & $2$ & $0$ &
$\mathcal I_{B_2}^{\mathrm{comp}}\simeq
 \mathcal I_{V_{10}}^{\mathrm{comp}}\simeq\MorseCat{3}^{\oplus2}$. \\
$1$ & $(B_1,V_6)$ & $2$ & $2$ & $1$ & $1$ & $0$ &
$\mathcal I_{B_1}^{\mathrm{comp}}\simeq\MorseCat{3}^{\oplus2}$,
$\mathcal I_{V_6}^{\mathrm{comp}}\simeq\MorseCat{3}$. \\
\end{tblr}
\end{center}

Thus the isolated compactified sectors agree for $d=2,3,4,5$.  For $d=3$,
$V_{14}$ has one open-torus Morse point and a second node created at the
boundary; for $d=5$, the additional boundary node of $V_{22}$ likewise
restores the multiplicity missing from its open torus.  For $d=1$, the
isolated sectors already have different numbers of Morse summands.

\subsection{Central-component counts}

Let $k(F)$ denote the number of irreducible components of the central fiber
of a minimal compactification of $W_F$.  For the standard rank-one models,
\begin{equation}
\label{eq:five_pair_component_formula}
k(F)=h^{1,2}(F)+1
\end{equation}
\cite{PrzyjalkowskiHodge}.  Write $h_B=h^{1,2}(B_d)$ and
$h_V=h^{1,2}(V_{4d+2})$.  Substitution gives
\begin{center}
\footnotesize
\begin{tblr}{
  width=0.86\textwidth,
  colspec={Q[c,wd=0.06\textwidth]Q[c,wd=0.17\textwidth]
    Q[c,wd=0.09\textwidth]Q[c,wd=0.10\textwidth]
    Q[c,wd=0.09\textwidth]Q[c,wd=0.10\textwidth]X[c]},
  colsep=2pt,
  row{1}={font=\bfseries},
  rowsep=2pt,
  hline{1,2,Z}={0.4pt},
}
$d$ & Pair & $h_B$ & $k(B_d)$ & $h_V$ & $k(V_{4d+2})$ & Status \\
$5$ & $(B_5,V_{22})$ & $0$ & $1$ & $0$ & $1$ & Equal. \\
$4$ & $(B_4,V_{18})$ & $2$ & $3$ & $2$ & $3$ & Equal. \\
$3$ & $(B_3,V_{14})$ & $5$ & $6$ & $5$ & $6$ & Equal. \\
$2$ & $(B_2,V_{10})$ & $10$ & $11$ & $10$ & $11$ & Equal. \\
$1$ & $(B_1,V_6)$ & $21$ & $22$ & $20$ & $21$ & Different by one. \\
\end{tblr}
\end{center}

Both available compactification counts therefore agree for $d=2,3,4,5$ and
fail to agree for $d=1$.  For $d=5$, both central fibers are irreducible.  For
$d=3$ and $d=2$, equality of the numbers $6$ and $11$, respectively, does not
determine the dual intersection complex, the transverse quadratic data, or
the categories at higher intersections.  Only for $d=4$ do the calculations
above determine the common incidence pattern: three components, three
rational double curves, and two triple points.

\subsection{The punctured double curves for the \texorpdfstring{$d=4$}{d=4} pair}

\begin{lemma}
\label{lem:punctured_double_curves}
For every $1\leq i<j\leq3$, one has
\[
R_{ij}^{\circ}
\cong
S_{ij}^{\circ}
\cong
\mathbb P^1\setminus\{0,\infty\}
\cong
\mathbb C^*,
\]
and the transverse potential on either curve has split normal form $uv$.
Consequently,
\[
\MBcat{R_{ij}^{\circ}}
\simeq
\Dtwo(R_{ij}^{\circ})
\simeq
\Dtwo(S_{ij}^{\circ})
\simeq
\MBcat{S_{ij}^{\circ}}.
\]
\end{lemma}

\begin{proof}
Each $R_{ij}$ and $S_{ij}$ is a smooth rational curve passing through the two
triple-point strata of the corresponding resolved central fiber.  Removing
those two points gives a twice-punctured projective line.  Away from the
triple points the two components meet transversely, so formal coordinates put
the potential in the form $uv$.  Relative Knörrer periodicity identifies the
matrix-factorization category with the two-periodic coherent category of the
supporting punctured curve.
\end{proof}

The lemma identifies the three open Morse--Bott pieces in
\eqref{eq:B4_resolved_localization} and
\eqref{eq:V18_resolved_localization}.  It does not identify their middle
terms or the complete enhanced Čech diagrams.  The precise additional
hypothesis is that there are dg equivalences
\[
\Phi_I:\mathfrak M^B_I\xrightarrow{\sim}\mathfrak M^V_I
\qquad
(\varnothing\neq I\subset\{\circ,X,Y,Z,T\})
\]
and coherent natural equivalences
\begin{equation}
\label{eq:d4_restriction_compatibility}
\Phi_J\circ\rho^B_{I,J}
\simeq
\rho^V_{I,J}\circ\Phi_I
\qquad(I\subset J),
\end{equation}
satisfying the cocycle and all higher coherence conditions.  Equivalently,
the $\Phi_I$ must assemble to a natural equivalence of semicosimplicial dg
diagrams
\[
\Phi^\bullet:
\check{\mathfrak M}^{B,\bullet}
\xrightarrow{\sim}
\check{\mathfrak M}^{V,\bullet}.
\]
Functoriality of the homotopy limit then gives an equivalence of dg
totalizations; applying $H^0$ and idempotent completion yields
\begin{equation}
\label{eq:d4_compatible_gluing}
\CcompB\simeq\CcompV.
\end{equation}
Thus the equality of the abstract punctured curves is necessary input for the
comparison, but it does not construct the equivalences $\Phi_I$, the overlap
homotopies, or their higher coherences.  In particular, the two outer terms
of a Verdier sequence are not being used to reconstruct its middle term.
For $d\neq4$, the component-count table alone does not supply analogous
double curves or gluing diagrams.

\section{Categorical interpretation and full-category equivalence}

\subsection{The genus-two categories for the \texorpdfstring{$d=4$}{d=4} pair}

The intersection of two quadrics $B_4\subset\mathbb P^5$ determines a pencil
of quadrics.  Its discriminant consists of six points of $\mathbb P^1$; the
double cover branched over them is a smooth genus-two curve, denoted
$C_{B_4}$.  The standard semiorthogonal decomposition is
\begin{equation}
\label{eq:B4_fano_sod}
\Db(B_4)
=
\left\langle
\Db(C_{B_4}),\mathcal O_{B_4},\mathcal O_{B_4}(H)
\right\rangle.
\end{equation}
This is the quadric-fibration description of
\cite{KuzQuadric,AuelBernardaraBolognesi}.

A smooth prime Fano threefold $V_{18}$ of genus $10$ is a codimension-two
linear section of the $G_2$ Grassmannian.  Homological projective duality
associates to it a double cover of a line branched at six points, hence another
smooth genus-two curve $C_{V_{18}}$.  Its derived category has the
semiorthogonal decomposition
\begin{equation}
\label{eq:V18_fano_sod}
\Db(V_{18})
=
\left\langle
\Db(C_{V_{18}}),\mathcal O_{V_{18}},\mathcal U_{V_{18}}^{\vee}
\right\rangle.
\end{equation}
The order of the two exceptional objects depends on the mutation convention;
the nontrivial component is $\Db(C_{V_{18}})$
\cite{KuzHyperplane,KuzFano}.

The relation between index-two degree-four Fano threefolds and index-one
genus-ten Fano threefolds was formulated categorically in \cite{KuzFano}.
The theorem of Kuznetsov--Shinder shows that the relevant nontrivial
components occur in a common deformation framework and corrects the scope of
the original conjecture \cite{KuzShinder}.  It does not identify the components
of an arbitrary pair.  In this note, a \emph{matched pair} means an additional
choice for which $C_{B_4}\cong C_{V_{18}}$; only under this hypothesis are the
two Fano-side nontrivial components immediately equivalent.

\subsection{What follows for the compactified matrix-factorization categories}

The preceding derived equivalences concern the coherent derived categories of
the Fano threefolds.  The categories $\CcompB$ and $\CcompV$ are instead
matrix-factorization categories on Landau--Ginzburg mirrors.  No reference
cited in this note states or proves either of the following equivalences for
the compactified central categories defined in
\eqref{eq:B4_cech_totalization} and \eqref{eq:V18_cech_totalization}:
\[
\CcompB\simeq\Dtwo(C_{B_4})
\qquad\text{or}\qquad
\CcompV\simeq\Dtwo(C_{V_{18}}).
\]
They are expected mirror-side identifications, not theorems imported from the
literature.  What Kuznetsov proves is the Fano-side pair of semiorthogonal
decompositions \eqref{eq:B4_fano_sod} and \eqref{eq:V18_fano_sod}; those
results do not by themselves identify the newly defined categories $\CcompB$
and $\CcompV$ with the periodic derived categories of the corresponding
curves.
The explicit resolution actually shows why a direct relative-Morse argument
cannot prove these formulas.  In both models the resolved central support is
not a smooth genus-two curve: it is a connected union of three rational curves
meeting at two triple points.  Consequently the proven descriptions are the
nonsplit localization sequences
\eqref{eq:B4_resolved_localization} and
\eqref{eq:V18_resolved_localization}, not identities with
$\Dtwo(C_{B_4})$ or $\Dtwo(C_{V_{18}})$.

Claiming either genus-two equivalence as a theorem would require one of the
following:
\begin{enumerate}
\item a direct equivalence between the connected rational-curve gluing category
in \eqref{eq:B4_resolved_localization} or
\eqref{eq:V18_resolved_localization} and the corresponding two-periodic
genus-two category;
\item a proved compactified Homological Mirror Symmetry equivalence that
identifies the central eigenvalue piece with the corresponding genus-two
category.
\end{enumerate}

The local calculation does prove the expected form away from the exceptional
strata: a split $uv$ normal form contributes the two-periodic coherent category
of the supporting smooth curve by Knörrer periodicity.  What remains is the
global extension over the crossing, the node, and the exceptional divisors
\cite{Knorrer,OrlovEquivalences,Hirano,LuntsSchnurer}.

It is therefore useful to state the expected identification separately.

\begin{remark}[Conditional genus-two identification]
Assume that compactified HMS, or a direct comparison of the two localization
sequences, identifies the connected formal gluing category of the three
rational curves with the genus-two component on the Fano side.  Then one
expects
\[
\mathcal C^{\mathrm{comp}}_{F,0}
\simeq
\Dtwo(C_F),
\qquad F\in\{B_4,V_{18}\}.
\]
For a Kuznetsov-matched pair with
$C_{B_4}\cong C_{V_{18}}$, this would imply
\[
\CcompB\simeq\CcompV.
\]
This implication is conditional; it is not a consequence of the five-chart
calculation by itself.
\end{remark}

\begin{corollary}[Equivalence for a matched pair]
\label{cor:matched_pair_equivalence}
Let $(B_4,V_{18})$ be a matched pair for which there is an isomorphism
\[
\phi:C_{B_4}\xrightarrow{\sim}C_{V_{18}}.
\]
Then pullback along $\phi$ gives
\[
\Db(C_{V_{18}})\simeq\Db(C_{B_4}),
\qquad
\Dtwo(C_{V_{18}})\simeq\Dtwo(C_{B_4}).
\]
If, in addition, the compactified mirror identifications
\[
\CcompB\simeq\Dtwo(C_{B_4}),
\qquad
\CcompV\simeq\Dtwo(C_{V_{18}})
\]
hold, then their composition yields
\begin{equation}
\label{eq:matched_compactified_equivalence}
\boxed{\CcompB\simeq\CcompV}.
\end{equation}
\end{corollary}

\begin{proof}
The isomorphism $\phi$ induces the first two equivalences.  Composing its
two-periodization with the two assumed compactified mirror equivalences gives
\[
\CcompB
\simeq\Dtwo(C_{B_4})
\simeq\Dtwo(C_{V_{18}})
\simeq\CcompV.
\]
\end{proof}

For $\mathbb Z$-graded derived categories of smooth curves of genus at least
two, an equivalence reconstructs the curve.  Thus
\[
\Db(C_{B_4})\simeq\Db(C_{V_{18}})
\quad\Longleftrightarrow\quad
C_{B_4}\cong C_{V_{18}}.
\]
After two-periodization, one should additionally require that an equivalence
admits a compatible $\mathbb Z$-graded lift before applying this reconstruction
statement.
The reconstruction implication for the $\mathbb Z$-graded categories is due
to Bondal--Orlov \cite{BondalOrlov}.

\subsection{Fano-side status for all five pairs}

Let $\operatorname{Ku}(B_d)$ and $\operatorname{Ku}(V_{4d+2})$ denote the
nontrivial components in the standard semiorthogonal decompositions.  Their
comparison is not uniform in $d$.
\begin{enumerate}[label=\textup{(\roman*)}]
\item For $d=5$,
\[
\operatorname{Ku}(B_5)
\simeq
\operatorname{Ku}(V_{22})
\simeq
\Db(\operatorname{Rep}K_3),
\]
where $K_3$ is the three-arrow Kronecker quiver \cite{KuzFano}.
\item For $d=4$ and $d=3$, equivalences hold for pairs lying in the
correspondences constructed by Kuznetsov; they are not statements about
arbitrary independently chosen pairs \cite{KuzFano}.
\item For $d=2$, the Kuznetsov components of a smooth quartic double solid
$B_2$ and a smooth Gushel--Mukai threefold $V_{10}$ are deformation
equivalent but not equivalent \cite{BayerPerry}.
\item For $d=1$, the corrected statement uses a nodal degree-one del Pezzo
threefold and a crepant categorical resolution of its nontrivial component,
rather than an arbitrary smooth pair $(B_1,V_6)$ \cite{KuzShinder}.
\end{enumerate}

Consequently, equality of isolated node counts or central-component counts
cannot by itself imply equivalence of the full compactified categories.  For
$d=3,4,5$, a matched Fano-side equivalence yields the expected mirror-side
equivalence only after compatible compactified-HMS identifications or a
direct equivalence of the global formal gluing diagrams.

\subsection{Case-by-case verification of the compactified equality}
\label{sec:case_by_case_compactified_equality}

There are two different equalities that must not be conflated.  Write
\[
\mathcal T_F^{\mathrm{comp}}
:=
\overline{\MF(Y_F,w_F)}
\]
for the total compactified matrix-factorization category, containing the
contributions from all critical values.  On the other hand, write
$\mathcal C_F^{\mathrm{comp}}$ for the distinguished compactified summand
expected, under compactified HMS, to correspond to
$\operatorname{Ku}(F)^{\mathrm{per}}$.  This distinguished summand need not
coincide with the category of the fiber labelled $0$ for an arbitrary
translation of the pencil parameter.

\begin{proposition}[Status of
$\mathcal C_{B_d}^{\mathrm{comp}}\simeq
\mathcal C_{V_{4d+2}}^{\mathrm{comp}}$]
\label{prop:case_by_case_compactified_equality}
The compactification calculations and the known Fano-side results give the
following case-by-case answer.
\begin{enumerate}[label=\textup{(\roman*)}]
\item For $d=5$, the total compactified categories are directly equivalent:
\begin{equation}
\label{eq:d5_total_compactified_equality}
\mathcal T_{B_5}^{\mathrm{comp}}
\simeq
\MorseCat{3}^{\oplus4}
\simeq
\mathcal T_{V_{22}}^{\mathrm{comp}}.
\end{equation}
The equality of the distinguished summands
$\mathcal C_{B_5}^{\mathrm{comp}}$ and
$\mathcal C_{V_{22}}^{\mathrm{comp}}$ follows from the equivalence of their
Kuznetsov components only after compatible compactified-HMS identifications.

\item For $d=4$ and $d=3$, let $(B_d,V_{4d+2})$ be a Kuznetsov-matched pair.
If
\[
\mathcal C_{B_d}^{\mathrm{comp}}
\simeq\operatorname{Ku}(B_d)^{\mathrm{per}},
\qquad
\mathcal C_{V_{4d+2}}^{\mathrm{comp}}
\simeq\operatorname{Ku}(V_{4d+2})^{\mathrm{per}},
\]
then
\begin{equation}
\label{eq:d3_d4_conditional_compactified_equality}
\mathcal C_{B_d}^{\mathrm{comp}}
\simeq
\mathcal C_{V_{4d+2}}^{\mathrm{comp}}.
\end{equation}
For $d=4$, the same conclusion follows directly if the complete enhanced
Čech diagrams are naturally equivalent, compatibly with all
restriction/incidence functors and their higher coherences, as required by
the criterion \eqref{eq:d4_compatible_gluing}.  Neither argument gives an
unconditional statement for an arbitrary independently chosen pair.

\item For $d=2$, an equivalence
\begin{equation}
\label{eq:d2_compactified_no_graded_lift}
\mathcal C_{B_2}^{\mathrm{comp}}
\simeq
\mathcal C_{V_{10}}^{\mathrm{comp}}
\end{equation}
cannot admit a compatible $\mathbb Z$-graded lift that matches graded
compactified HMS.  Such a lift would imply an equivalence between the
Kuznetsov components of a smooth quartic double solid and a smooth
Gushel--Mukai threefold, contradicting \cite{BayerPerry}.  In the purely
ungraded two-periodic setting, the present calculation neither proves nor
disproves \eqref{eq:d2_compactified_no_graded_lift}.

\item For $d=1$, the expression
$\mathcal C_{B_1}^{\mathrm{comp}}\simeq
\mathcal C_{V_6}^{\mathrm{comp}}$ for an arbitrary smooth pair is not the
statement furnished by the corrected Fano correspondence.  The appropriate
comparison replaces the smooth $B_1$ by a nodal degree-one del Pezzo
threefold and replaces its Kuznetsov component by a crepant categorical
resolution \cite{KuzShinder}.  A compactified mirror equivalence for this
corrected pair is not proved here.
\end{enumerate}
Consequently, the formula
\[
\mathcal C_{B_d}^{\mathrm{comp}}
\simeq
\mathcal C_{V_{4d+2}}^{\mathrm{comp}}
\qquad(1\leq d\leq5)
\]
is not a theorem.  It is conditionally valid for matched pairs when
$d=3,4,5$, has no compatible graded-HMS lift for smooth pairs when $d=2$,
and must be reformulated when $d=1$.
\end{proposition}

\begin{proof}
For $d=5$, equation \eqref{eq:d5_total_compactified_equality} is the
combination of \eqref{eq:B5_compact_final_decomposition} and
\eqref{eq:V22_compact_final_decomposition}.  The distinguished-summand
statement follows by two-periodizing
$\operatorname{Ku}(B_5)\simeq\operatorname{Ku}(V_{22})$ and composing with
the two assumed mirror identifications.

For $d=3,4$, two-periodize the equivalence of the Kuznetsov components of a
matched pair and compose it with the two assumed compactified-HMS
equivalences.  The alternative $d=4$ argument is precisely the homotopy-gluing
criterion \eqref{eq:d4_compatible_gluing}.

For $d=2$, a compatible graded lift of
\eqref{eq:d2_compactified_no_graded_lift}, composed with graded compactified
HMS, would give
$\operatorname{Ku}(B_2)\simeq\operatorname{Ku}(V_{10})$, contradicting the
nonequivalence theorem of \cite{BayerPerry}.  Forgetting to the two-periodic
category loses this obstruction, so no ungraded nonequivalence is claimed.
The final statement is the corrected degeneration theorem of
\cite{KuzShinder}.
\end{proof}

\begin{remark}[The fiber labelled zero for $d=5$]
\label{rem:d5_zero_fiber_warning}
For the explicit parameter normalizations used above,
\[
\mathcal C_{B_5,0}^{\mathrm{comp}}=0,
\qquad
\mathcal C_{V_{22},0}^{\mathrm{comp}}\simeq\MorseCat{3},
\]
because the $B_5$ fiber over $0$ is smooth whereas the fourth $V_{22}$ node
lies over $0$.  Thus the categories of the fibers carrying the literal label
$0$ are not equivalent.  Equation \eqref{eq:d5_total_compactified_equality}
compares the total categories, while the distinguished-summand comparison
requires a matching of critical values or compactified HMS.
\end{remark}

\begin{theorem}[Compactified comparison for the five pairs]
\label{thm:five_pair_conclusion}
For the standard compactifications of $(B_d,V_{4d+2})$, $1\leq d\leq5$, the
following statements hold.
\begin{enumerate}[label=\textup{(\alph*)}]
\item The five-chart calculation for $d=4$ gives
\[
\overline{\MF(Y_{B_4},w_{B_4})}
\simeq\CcompB\oplus\MorseCat{3}^{\oplus2},
\]
and
\[
\overline{\MF(Y_{V_{18}},w_{V_{18}})}
\simeq\CcompV\oplus\MorseCat{3}^{\oplus2}.
\]
The central categories are the Čech totalizations in
\eqref{eq:B4_cech_totalization} and \eqref{eq:V18_cech_totalization}.
\item The explicit $d=2$ calculation gives
\[
\overline{\MF(Y_{B_2},w_{B_2})}
\simeq\CcompBtwo\oplus\MorseCat{3}^{\oplus2},
\qquad
\overline{\MF(Y_{V_{10}},w_{V_{10}})}
\simeq\CcompVten\oplus\MorseCat{3}^{\oplus2}.
\]
The raw central supports are, respectively, a line together with an isolated
point and a double smooth quadric.  Thus the isolated sectors agree, but the
calculation does not identify $\CcompBtwo$ with $\CcompVten$.
\item The explicit $d=3$ calculation gives
\[
\overline{\MF(Y_{B_3},w_{B_3})}
\simeq\CcompBthree\oplus\MorseCat{3}^{\oplus2},
\qquad
\overline{\MF(Y_{V_{14}},w_{V_{14}})}
\simeq\CcompVfourteen\oplus\MorseCat{3}^{\oplus2}.
\]
For $B_3$, both isolated ordinary double points are visible on the torus.
For $V_{14}$, one occurs on the torus at critical value $27$, whereas the
second appears over the boundary base curve in the fiber of critical value
$-1$.  The isolated sectors and the resolved central-component counts agree,
but the calculation does not identify $\CcompBthree$ with
$\CcompVfourteen$.
\item The compactification calculation on the index-two side is complete for
all five values of $d$.  In particular,
\[
\overline{\MF(Y_{B_1},w_{B_1})}
\simeq\CcompBone\oplus\MorseCat{3}^{\oplus2},
\qquad
\overline{\MF(Y_{B_5},w_{B_5})}
\simeq\MorseCat{3}^{\oplus4}.
\]
Their central-component counts are
\[
\bigl(k(B_1),k(B_2),k(B_3),k(B_4),k(B_5)\bigr)
=(22,11,6,3,1).
\]
For $B_5$ the resolved pencil is Lefschetz and its fiber over $0$ is smooth.
\item The two extreme index-one compactifications give
\[
\overline{\MF(Y_{V_6},w_{V_6})}
\simeq\CcompVsix\oplus\MorseCat{3},
\qquad
\overline{\MF(Y_{V_{22}},w_{V_{22}})}
\simeq\MorseCat{3}^{\oplus4}.
\]
For $V_6$, the unique noncentral node has critical value $108$ and the
resolved central fiber has $21$ components.  For $V_{22}$, three nodes lie
away from the base locus and the fourth is created by blowing up the line
$T=Y+Z=0$; the pencil is Lefschetz.
\item The isolated compactified categories satisfy
\[
\mathcal I_{B_d}^{\mathrm{comp}}
\simeq
\mathcal I_{V_{4d+2}}^{\mathrm{comp}}
\quad\Longleftrightarrow\quad
d\in\{2,3,4,5\}.
\]
For $d=5$, the equality uses the additional boundary ordinary double point of
$V_{22}$.
\item The central-component counts satisfy
\[
k(B_d)=k(V_{4d+2})
\quad\Longleftrightarrow\quad
d\in\{2,3,4,5\}.
\]
For $d=1$, $k(B_1)=22$ and $k(V_6)=21$, and the isolated sectors also have
different multiplicities.
\item For $d=4$, the three punctured double-curve categories are pairwise
equivalent by Lemma~\ref{lem:punctured_double_curves}.  If the complete
enhanced Čech diagrams are naturally equivalent, compatibly with all
restriction/incidence functors and their higher coherences, then
$\CcompB\simeq\CcompV$.
\item For $d\in\{3,4,5\}$ and a Kuznetsov-matched pair, compatible
compactified-HMS identifications imply
\[
\mathcal C^{\mathrm{comp}}_{B_d}
\simeq
\mathcal C^{\mathrm{comp}}_{V_{4d+2}}.
\]
This is conditional: the numerical compactification data do not construct
the required equivalence of gluing diagrams.  For $d=2$ the corresponding
smooth Fano-side components are not equivalent, and for $d=1$ the statement
must be reformulated using the nodal degree-one del Pezzo category.
\end{enumerate}
Therefore the compactification calculation proves equality of the isolated
sectors for exactly four pairs.  Equivalence of the full compactified
categories requires the additional global gluing or compactified-HMS input
specified above.
\end{theorem}

\begin{proof}
Part~\textup{(a)} is the combination of
\eqref{eq:B4_final_decomposition} and
\eqref{eq:V18_final_decomposition}.  Part~\textup{(b)} is
Proposition~\ref{prop:d2_compactified_comparison} together with
\eqref{eq:B2_compact_final_decomposition} and
\eqref{eq:V10_compact_final_decomposition}.  Part~\textup{(c)} is
Proposition~\ref{prop:d3_compactified_comparison} together with
\eqref{eq:B3_compact_final_decomposition} and
\eqref{eq:V14_compact_final_decomposition}.  Part~\textup{(d)} is
Proposition~\ref{prop:Bd_compactification_census}.  Part~\textup{(e)} is
Proposition~\ref{prop:V6_V22_compactification_census}.  Part~\textup{(f)}
follows from \eqref{eq:five_pair_isolated_category} and the
ordinary-double-point table.  Part~\textup{(g)} follows from
\eqref{eq:five_pair_component_formula} and the Hodge-number table.
Part~\textup{(h)} is Lemma~\ref{lem:punctured_double_curves} together with the
compatible-gluing implication \eqref{eq:d4_compatible_gluing}.  Part~\textup{(i)}
follows by two-periodizing the matched Fano-side equivalence and composing it
with the two assumed compactified mirror identifications.  The qualifications
for $d=1,2$ are the results recalled in the preceding subsection.
\end{proof}

\begin{thebibliography}{99}

\bibitem{Minkowski}
M.~Akhtar, T.~Coates, S.~Galkin, and A.~M.~Kasprzyk,
\emph{Minkowski polynomials and mutations},
SIGMA \textbf{8} (2012), Paper 094.

\bibitem{AuelBernardaraBolognesi}
A.~Auel, M.~Bernardara, and M.~Bolognesi,
\emph{Fibrations in complete intersections of quadrics, Clifford algebras,
derived categories, and rationality problems},
J. Math. Pures Appl. \textbf{102} (2014), 249--291.

\bibitem{BondalOrlov}
A.~Bondal and D.~Orlov,
\emph{Reconstruction of a variety from the derived category and groups of
autoequivalences},
Compos. Math. \textbf{125} (2001), 327--344.

\bibitem{CCGK}
T.~Coates, A.~Corti, S.~Galkin, and A.~Kasprzyk,
\emph{Quantum periods for $3$-dimensional Fano manifolds},
Geom. Topol. \textbf{20} (2016), 103--256.

\bibitem{Dyckerhoff}
T.~Dyckerhoff,
\emph{Compact generators in categories of matrix factorizations},
Duke Math. J. \textbf{159} (2011), 223--274.

\bibitem{Eisenbud}
D.~Eisenbud,
\emph{Homological algebra on a complete intersection, with an application to
group representations},
Trans. Amer. Math. Soc. \textbf{260} (1980), 35--64.

\bibitem{Hirano}
Y.~Hirano,
\emph{Derived Kn\"orrer periodicity and Orlov's theorem for gauged
Landau--Ginzburg models},
Compos. Math. \textbf{153} (2017), 973--1007.

\bibitem{Knorrer}
H.~Kn\"orrer,
\emph{Cohen--Macaulay modules on hypersurface singularities I},
Invent. Math. \textbf{88} (1987), 153--164.

\bibitem{KuzHyperplane}
A.~Kuznetsov,
\emph{Hyperplane sections and derived categories},
Izv. Math. \textbf{70} (2006), 447--547.

\bibitem{KuzQuadric}
A.~Kuznetsov,
\emph{Derived categories of quadric fibrations and intersections of quadrics},
Adv. Math. \textbf{218} (2008), 1340--1369.

\bibitem{KuzFano}
A.~Kuznetsov,
\emph{Derived categories of Fano threefolds},
Proc. Steklov Inst. Math. \textbf{264} (2009), 110--122;
arXiv:0809.0225.

\bibitem{BayerPerry}
A.~Bayer and A.~Perry,
\emph{Kuznetsov's Fano threefold conjecture via K3 categories and enhanced
group actions},
J. Reine Angew. Math. \textbf{800} (2023), 107--153.

\bibitem{KuzShinder}
A.~Kuznetsov and E.~Shinder,
\emph{Derived categories of Fano threefolds and degenerations},
Invent. Math. \textbf{239} (2025), 377--430.

\bibitem{LuntsSchnurer}
V.~A.~Lunts and O.~M.~Schn\"urer,
\emph{Matrix factorizations and semi-orthogonal decompositions for
blowing-ups},
J. Noncommut. Geom. \textbf{10} (2016), 907--979.

\bibitem{OrlovDbranes}
D.~Orlov,
\emph{Triangulated categories of singularities and D-branes in
Landau--Ginzburg models},
Proc. Steklov Inst. Math. \textbf{246} (2004), 227--248.

\bibitem{OrlovEquivalences}
D.~Orlov,
\emph{Triangulated categories of singularities and equivalences between
Landau--Ginzburg models},
Sb. Math. \textbf{197} (2006), 1827--1840.

\bibitem{OrlovFormal}
D.~Orlov,
\emph{Formal completions and idempotent completions of triangulated
categories of singularities},
Adv. Math. \textbf{226} (2011), 206--217.

\bibitem{OrlovNonaffine}
D.~Orlov,
\emph{Matrix factorizations for nonaffine Landau--Ginzburg models},
Math. Ann. \textbf{353} (2012), 95--108.

\bibitem{Preygel}
A.~Preygel,
\emph{Thom--Sebastiani and duality for matrix factorizations},
arXiv:1101.5834.

\bibitem{PrinceMutation}
T.~Prince,
\emph{Mutation equivalence of toric Landau--Ginzburg models},
arXiv:2006.01477.

\bibitem{PrzyjalkowskiWeak}
V.~Przyjalkowski,
\emph{Weak Landau--Ginzburg models for smooth Fano threefolds},
Izv. Math. \textbf{77} (2013), 772--794.

\bibitem{PrzyjalkowskiHodge}
V.~Przyjalkowski,
\emph{Hodge numbers of Fano threefolds via Landau--Ginzburg models},
arXiv:0911.5428 (2009).

\bibitem{PrzyjalkowskiCY}
V.~Przyjalkowski,
\emph{Calabi--Yau compactifications of toric Landau--Ginzburg models for
smooth Fano threefolds},
Sb. Math. \textbf{208} (2017), 992--1013.

\bibitem{PrzyjalkowskiToric}
V.~Przyjalkowski,
\emph{Toric Landau--Ginzburg models},
Russian Math. Surveys \textbf{73} (2018), 1033--1118.

\bibitem{Przyjalkowski2026}
V.~Przyjalkowski,
\emph{Landau--Ginzburg models for Fano threefolds of Picard rank one and
exceptional collections},
arXiv:2601.16497 (2026).

\end{thebibliography}

\end{document}